\documentclass[10pt,a4paper]{article}
\usepackage{../../styles/neet_style}

\pagestyle{fancy}
\fancyhf{}
\fancyhead[L]{\small \textbf{NEET UG 2027: Inorganic Chemistry Master Compendium}}
\fancyhead[R]{\small Satya Prakash}
\fancyfoot[L]{\small \textbf{NEET UG 2027 Prep Series}}
\fancyfoot[C]{\thepage}
\fancyfoot[R]{\small \textbf{Inorganic Chemistry Chapter Notes}}

\begin{document}

% Front Banner
\begin{center}
  {\Huge \textbf{NEET UG 2027: Master Chapter Compendium}}\\[6pt]
  {\Large \textbf{Inorganic Chemistry: In-Depth Structural Theory \& Comparative Principles}}\\[4pt]
  {\small \textbf{Prepared, Compiled \& Authored by Satya Prakash} $\cdot$ Strictly Calibrated to Revised NMC Syllabus}
  \vspace{4pt}
  \hrule height 1.5pt
\end{center}

% =========================================================================
% SECTION 1: NMC SYLLABUS DIRECTIVE & STRICT EXCLUSIONS
% =========================================================================
\section{NMC 2027 Regulatory Calibration \& Scope Demarcation}

\begin{trapbox}[Zero-Tolerance NMC Syllabus Alert: Deleted Inorganic Chemistry Chapters]
Under the revised National Medical Commission (NMC) NEET UG 2027 regulatory curriculum, the following legacy chapters have been \textbf{strictly deleted and removed from the examination}:
\begin{enumerate}
  \item \textbf{Hydrogen:} Position in periodic table, isotopes, preparation and properties of dihydrogen, hydrides (ionic, covalent, interstitial), physical and chemical properties of water, hard/soft water, heavy water ($\mathrm{D_2O}$), and hydrogen peroxide ($\mathrm{H_2O_2}$) are \textbf{completely eliminated}.
  \item \textbf{$s$-Block Elements (Alkali \& Alkaline Earth Metals):} Group 1 and Group 2 chemistry, electronic configurations, diagonal relationships, flame coloration, and industrial compounds ($\mathrm{Na_2CO_3, NaCl, NaOH, NaHCO_3, CaO, CaCO_3, CaSO_4\cdot \frac{1}{2}H_2O}$, cement) are \textbf{completely eliminated}.
  \item \textbf{General Principles \& Processes of Isolation of Elements (Metallurgy):} Ores, mineral dressing, roasting, calcination, smelting, Ellingham diagrams, blast furnace iron extraction, copper reverberatory smelting, Hall-Héroult aluminum extraction, and refining techniques (Mond process, van Arkel method, zone refining) are \textbf{completely eliminated}.
  \item \textbf{Environmental Chemistry:} Tropospheric/stratospheric pollution, classical and photochemical smog, ozone depletion, and green chemistry are \textbf{strictly eliminated}.
\end{enumerate}
\textbf{Study Directive:} Do not invest study hours in these eliminated topics. All cognitive focus must remain on the rigorously tested Class 11 and Class 12 inorganic chemistry compendium units detailed below.
\end{trapbox}

% =========================================================================
% SECTION 2: PERIODICITY & PERIODIC ANOMALIES
% =========================================================================
\section{Classification of Elements \& Periodicity in Properties (Unit 9)}

\subsection{Modern Periodic Law \& Electronic Architecture}
Moseley's Law (1913) demonstrated that the square root of the characteristic X-ray frequency ($\sqrt{\nu}$) emitted by an element is directly proportional to its atomic number $Z$:
\begin{equation}
  \sqrt{\nu} = a(Z - b)
\end{equation}
Thus, the physical and chemical properties of the elements are periodic functions of their \textbf{atomic numbers}, not their atomic masses. The long form of the periodic table consists of 7 horizontal periods and 18 vertical groups, partitioned into $s, p, d,$ and $f$ blocks based on the differentiating electron orbital.

\subsection{Effective Nuclear Charge ($Z_{\text{eff}}$) \& Slater's Rules}
Electrons in outer shells are electrostatically shielded from the full nuclear charge $+Ze$ by intervening inner-shell electron density:
\begin{equation}
  Z_{\text{eff}} = Z - \sigma
\end{equation}
where $\sigma$ is the screening or shielding constant.
\begin{itemize}
  \item \textbf{Shielding Efficiency Hierarchy:} $s > p > d > f$.
  Spherical $s$-orbitals penetrate closest to the nucleus, providing maximal screening. Diffuse $d$- and $f$-orbitals have poor radial penetration and provide \textbf{extraordinarily poor shielding}.
  \item \textbf{Slater's Rules Formulation for $ns, np$ Valence Electrons:}
  \begin{enumerate}
    \item Electrons in shells higher than the considered electron contribute $0.00$.
    \item Other electrons in the identical $(ns, np)$ group contribute $0.35$ each (or $0.30$ if $1s$).
    \item Each electron in the $(n - 1)$ shell contributes $0.85$.
    \item Each electron in the $(n - 2)$ or lower inner shells contributes $1.00$.
  \end{enumerate}
  \item \textbf{For $nd, nf$ Valence Electrons:}
  \begin{enumerate}
    \item Other electrons in the identical $(nd, nf)$ group contribute $0.35$ each.
    \item All electrons in all lower underlying shells contribute $1.00$ each.
  \end{enumerate}
\end{itemize}

\begin{usecasebox}[High-Yield Problem Archetype: Step-by-Step Slater's $Z_{\text{eff}}$ Calculation]
\textbf{Archetype Problem:} Calculate the effective nuclear charge ($Z_{\text{eff}}$) experienced by: (a) a $4s$ valence electron in Zinc ($Z=30$), and (b) a $3d$ electron in Zinc ($Z=30$).
\begin{enumerate}
  \item \textbf{Electronic Configuration Grouping:}
  \begin{equation*}
    (1s)^2 \ (2s, 2p)^8 \ (3s, 3p)^8 \ (3d)^{10} \ (4s)^2
  \end{equation*}
  \item \textbf{For the $4s$ Valence Electron ($n=4$):}
  \begin{itemize}
    \item $1$ other electron in same group ($4s$): $1 \times 0.35 = 0.35$.
    \item $(n-1)$ shell (both $3s, 3p$ and $3d$, total 18 electrons): $18 \times 0.85 = 15.30$.
    \item $(n-2)$ and lower shells ($1s, 2s, 2p$, total 10 electrons): $10 \times 1.00 = 10.00$.
    \item $\sigma_{4s} = 0.35 + 15.30 + 10.00 = 25.65 \implies Z_{\text{eff}} = 30 - 25.65 = \mathbf{4.35}$.
  \end{itemize}
  \item \textbf{For a $3d$ Electron ($n=3$):}
  \begin{itemize}
    \item $4s$ electrons lie in a higher shell $\implies$ contribute $0.00$.
    \item $9$ other electrons in the same $(3d)$ group: $9 \times 0.35 = 3.15$.
    \item All 18 electrons in lower shells: $18 \times 1.00 = 18.00$.
    \item $\sigma_{3d} = 3.15 + 18.00 = 21.15 \implies Z_{\text{eff}} = 30 - 21.15 = \mathbf{8.85}$.
  \end{itemize}
  \item \textbf{Conclusion:} $3d$ electrons experience more than double the effective nuclear attraction ($8.85$ vs $4.35$) compared to $4s$ electrons. Hence, upon ionization, the $4s$ electrons are expelled first: $\mathrm{Zn \to Zn^{2+} + 2e^-}$ yields $[\mathrm{Ar}] 3d^{10}$, not $[\mathrm{Ar}] 3d^8 4s^2$!
\end{enumerate}
\end{usecasebox}

\subsection{Atomic and Ionic Radii Trends \& Contractions}
\begin{itemize}
  \item \textbf{Radius Hierarchy:} For the same element: $r_{\text{cation}} < r_{\text{atom}} < r_{\text{anion}}$.
  Loss of electrons increases $Z_{\text{eff}} = Z/e$, pulling electrons closer. Gain of electrons increases interelectronic repulsion, expanding the electron cloud.
  \item \textbf{Isoelectronic Series Size Rule:} For ions sharing the identical total electron count:
  \begin{equation}
    r_{\text{isoelectronic}} \propto \frac{1}{Z} \implies \mathrm{C^{4-} > N^{3-} > O^{2-} > F^- > Na^+ > Mg^{2+} > Al^{3+} > Si^{4+}} \quad (10e^-)
  \end{equation}
  \item \textbf{$3d$-Transition Metal Contraction (d-Block Contraction):}
  Filling of the ten $3d$ orbitals across the first transition series precedes Gallium ($Z=31$: $[\mathrm{Ar}] 3d^{10} 4s^2 4p^1$). Because $3d$ electrons shield poorly, $Z_{\text{eff}}$ on the valence $4p$ electron is abnormally high. Consequently:
  \begin{equation}
    r(\mathrm{Ga}) = 135\,\si{pm} \le r(\mathrm{Al}) = 143\,\si{pm}
  \end{equation}
  \item \textbf{Lanthanoid Contraction ($4f$-Contraction):}
  Filling of the fourteen diffuse $4f$ orbitals between Lanthanum ($Z=57$) and Hafnium ($Z=72$) results in extremely poor screening of the increasing nuclear charge ($+14$ units).
  \textbf{Consequences:}
  \begin{enumerate}
    \item Nearly identical covalent and ionic radii for $4d$ and $5d$ congeneric pairs:
    \begin{equation}
      r(\mathrm{Zr}) \approx r(\mathrm{Hf}) \approx 159-160\,\si{pm}, \qquad r(\mathrm{Nb}) \approx r(\mathrm{Ta}) \approx 146\,\si{pm}, \qquad r(\mathrm{Mo}) \approx r(\mathrm{W}) \approx 139\,\si{pm}
    \end{equation}
    \item High density of $5d$ elements (roughly double their $4d$ counterparts; e.g., Osmium $\rho \approx 22.6\,\si{g.cm^{-3}}$).
    \item Regular decrease in basicity of lanthanoid hydroxides: $\mathrm{La(OH)_3 > Ce(OH)_3 > \dots > Lu(OH)_3}$ (covalent character increases as cation size contracts).
  \end{enumerate}
\end{itemize}

\subsection{Ionization Enthalpy ($\Delta_i H$) \& High-Yield Anomalies}
Ionization Enthalpy is the minimum energy required to remove the most loosely bound electron from an isolated gaseous atom in its ground state:
\begin{equation}
  \mathrm{X(g)} \xrightarrow{\text{IE}_1} \mathrm{X^+(g)} + e^-, \qquad \mathrm{X^+(g)} \xrightarrow{\text{IE}_2} \mathrm{X^{2+}(g)} + e^-
\end{equation}
Successive values always increase: $\text{IE}_1 < \text{IE}_2 < \text{IE}_3 < \dots$ because removing an electron from an increasingly positive ion encounters greater electrostatic attraction.

\begin{table}[htbp]
\centering
\small
\renewcommand{\arraystretch}{1.3}
\begin{tabularx}{\textwidth}{p{3.2cm} p{4.2cm} p{8.8cm}}
\toprule
\textbf{Element Pair} & \textbf{Observed Trend} & \textbf{Rigorous Theoretical Origin \& NCERT Rationale} \\
\midrule
\textbf{Beryllium vs Boron} & $\text{IE}_1(\mathrm{Be}) = 899\,\si{kJ.mol^{-1}}$ \newline $> \text{IE}_1(\mathrm{B}) = 801\,\si{kJ.mol^{-1}}$ & $\mathrm{Be}$ ($1s^2 2s^2$) loses a strongly penetrating $2s$ electron from a fully-filled stable subshell. $\mathrm{B}$ ($1s^2 2s^2 2p^1$) loses a less penetrating, well-shielded $2p$ electron of higher orbital energy. \\
\midrule
\textbf{Nitrogen vs Oxygen} & $\text{IE}_1(\mathrm{N}) = 1402\,\si{kJ.mol^{-1}}$ \newline $> \text{IE}_1(\mathrm{O}) = 1314\,\si{kJ.mol^{-1}}$ & $\mathrm{N}$ ($1s^2 2s^2 2p_x^1 2p_y^1 2p_z^1$) possesses an extra-stable, exactly half-filled $2p^3$ subshell with maximum exchange energy. $\mathrm{O}$ ($1s^2 2s^2 2p_x^2 2p_y^1 2p_z^1$) contains two electrons paired in the same $2p_x$ orbital; mutual Coulombic repulsion destabilizes the electron, making it easier to expel. \\
\midrule
\textbf{Magnesium vs Aluminum} & $\text{IE}_1(\mathrm{Mg}) = 737\,\si{kJ.mol^{-1}}$ \newline $> \text{IE}_1(\mathrm{Al}) = 577\,\si{kJ.mol^{-1}}$ & Exactly analogous to $\mathrm{Be/B}$: $\mathrm{Mg}$ has fully-filled $3s^2$ subshell; $\mathrm{Al}$ has solitary $3p^1$ electron outside neon core. \\
\midrule
\textbf{Phosphorus vs Sulfur} & $\text{IE}_1(\mathrm{P}) = 1012\,\si{kJ.mol^{-1}}$ \newline $> \text{IE}_1(\mathrm{S}) = 1000\,\si{kJ.mol^{-1}}$ & Exactly analogous to $\mathrm{N/O}$: $\mathrm{P}$ has half-filled $3p^3$ subshell; $\mathrm{S}$ has paired $3p_x^2 3p_y^1 3p_z^1$ experiencing spin-pairing repulsion. \\
\midrule
\textbf{Thallium vs Indium} & $\text{IE}_1(\mathrm{Tl}) = 589\,\si{kJ.mol^{-1}}$ \newline $> \text{IE}_1(\mathrm{In}) = 558\,\si{kJ.mol^{-1}}$ & Intervening $4f^{14}$ subshell in Thallium provides poor shielding (Lanthanoid contraction), causing high $Z_{\text{eff}}$ on outer $6s^2 6p^1$ electrons. \\
\midrule
\textbf{Lead vs Tin} & $\text{IE}_1(\mathrm{Pb}) = 715\,\si{kJ.mol^{-1}}$ \newline $> \text{IE}_1(\mathrm{Sn}) = 708\,\si{kJ.mol^{-1}}$ & Identical origin: $4f^{14}$ Lanthanoid contraction raises $Z_{\text{eff}}$ in Lead, arresting the normal downward decrease in IE. \\
\bottomrule
\end{tabularx}
\caption{Systematic Catalogue of High-Yield Ionization Enthalpy Anomalies.}
\end{table}

\subsection{Electron Gain Enthalpy ($\Delta_{\text{eg}} H$) Trends \& Traps}
Electron gain enthalpy is the enthalpy change accompanying the addition of an electron to an isolated gaseous atom to form a uninegative gaseous anion:
\begin{equation}
  \mathrm{X(g)} + e^- \longrightarrow \mathrm{X^-(g)}, \qquad \Delta_{\text{eg}} H
\end{equation}
Most elements release energy upon gaining an electron ($\Delta_{\text{eg}} H < 0$). However, second electron gain enthalpies are \textbf{strictly positive (endothermic)} for all elements (e.g., $\mathrm{O^-(g)} + e^- \to \mathrm{O^{2-}(g)}$, $\Delta_{\text{eg}} H_2 = +780\,\si{kJ.mol^{-1}}$) due to overpowering electrostatic repulsion between the incoming electron and the existing negative anion.

\begin{trapbox}[Cardinal Periodic Trap: Fluorine vs Chlorine \& Oxygen vs Sulfur]
\begin{itemize}
  \item \textbf{Halogen Group 17 Order:}
  \begin{equation}
    |\Delta_{\text{eg}} H|: \mathbf{\mathrm{Cl > F > Br > I}} \implies \Delta_{\text{eg}} H: \mathrm{Cl (-349) < F (-328) < Br (-325) < I (-295)\,\si{kJ.mol^{-1}}}
  \end{equation}
  \textbf{Reason:} Fluorine has an extraordinarily compact $2p$ valence shell ($n=2$). The incoming electron encounters intense interelectronic repulsion from the 7 tightly packed valence electrons. In Chlorine, the larger, more diffuse $3p$ shell accommodates the incoming electron with substantially less Coulombic repulsion.
  \item \textbf{Chalcogen Group 16 Order:}
  \begin{equation}
    |\Delta_{\text{eg}} H|: \mathbf{\mathrm{S > Se > Te > Po > O}} \implies \text{Oxygen has the \textbf{LEAST negative} value in Group 16!}
  \end{equation}
  Oxygen ($\Delta_{\text{eg}} H = -141\,\si{kJ.mol^{-1}}$) suffers severe electron-electron repulsion in its tiny $2p$ shell, whereas Sulfur ($\Delta_{\text{eg}} H = -200\,\si{kJ.mol^{-1}}$) accommodates the electron far more readily.
  \item \textbf{Positive / Near-Zero $\Delta_{\text{eg}} H$ Values:}
  \begin{itemize}
    \item \textbf{Noble Gases ($\mathrm{He, Ne, Ar, Kr, Xe}$):} Strongly positive $\Delta_{\text{eg}} H > 0$ because the incoming electron must enter a higher unpopulated principal shell.
    \item \textbf{Alkaline Earth Metals ($\mathrm{Be, Mg}$):} $\Delta_{\text{eg}} H \approx 0$ due to fully-filled stable $ns^2$ subshells.
    \item \textbf{Nitrogen ($\mathrm{N}$):} $\Delta_{\text{eg}} H \approx +31\,\si{kJ.mol^{-1}}$ due to extra-stable half-filled $2p^3$ configuration.
  \end{itemize}
\end{itemize}
\end{trapbox}

\subsection{Electronegativity \& The Inert Pair Effect}
\begin{itemize}
  \item \textbf{Electronegativity ($\chi$):} Qualitative tendency of an atom in a covalent molecule to attract shared bonding electrons toward itself.
  \begin{itemize}
    \item \textbf{Pauling Scale:} Based on bond dissociation energy discrepancies:
    \begin{equation}
      \Delta = E_{A-B} - \sqrt{E_{A-A} \times E_{B-B}} \implies |\chi_A - \chi_B| = 0.208\sqrt{\Delta\text{ (in kcal.mol}^{-1}\text{)}} = 0.102\sqrt{\Delta\text{ (in kJ.mol}^{-1}\text{)}}
    \end{equation}
    Pauling assigned Fluorine the benchmark value $\chi(\mathrm{F}) = 4.0$. Absolute hierarchy:
    \begin{equation}
      \mathbf{\mathrm{F (4.0) > O (3.5) > N (3.0) \approx Cl (3.0) > Br (2.8) > I (2.5) \approx C (2.5) \approx S (2.5) > H (2.1)}}
    \end{equation}
    \item \textbf{Mulliken Scale:} Arithmetic mean of ionization enthalpy and electron affinity:
    \begin{equation}
      \chi_{\text{Mulliken}} = \frac{\text{IE} + \text{EA}}{2}\text{ (in eV)}, \qquad \chi_{\text{Pauling}} \approx \frac{\chi_{\text{Mulliken}}}{2.8} = \frac{\text{IE} + \text{EA}}{5.6}
    \end{equation}
    \item \textbf{Hybridization Influence:} Electronegativity increases with increasing $\% s$-character:
    \begin{equation}
      sp\text{ (50\% } s\text{)} > sp^2\text{ (33.3\% } s\text{)} > sp^3\text{ (25\% } s\text{)}
    \end{equation}
  \end{itemize}
  \item \textbf{The Inert Pair Effect:}
  The reluctance of the valence $ns^2$ electron pair to participate in chemical bonding down Groups 13, 14, and 15.
  \textbf{Physical Origin:}
  \begin{enumerate}
    \item Intervening $(n-1)d^{10}$ and $(n-2)f^{14}$ electrons shield the nucleus poorly, subjecting the spherical $ns$ orbital to an abnormally high effective nuclear charge $Z_{\text{eff}}$.
    \item Relativistic contraction pulls $ns$ electrons closer to the nucleus, significantly increasing the energy required to unpair and promote an $s$-electron to the $p$-subshell.
    \item Bond energies decrease down the group due to larger orbital size; the energy released by forming two additional covalent bonds is insufficient to compensate for the promotion energy of the $ns^2$ electrons.
  \end{enumerate}
  \textbf{Consequences on Oxidation State Stabilities:}
  \begin{itemize}
    \item \textbf{Group 13:} Stability of $+1$ state increases: $\mathrm{B^{3+}} > \mathrm{Al^{3+}} > \mathrm{Ga^{3+}} > \mathrm{In^{3+}} > \mathrm{Tl^{3+}}$, whereas $\mathbf{\mathrm{Tl^+}} > \mathbf{\mathrm{In^+}} > \mathbf{\mathrm{Ga^+}} > \mathbf{\mathrm{Al^+}}$. $\mathrm{TlCl}$ is stable; $\mathrm{Tl^{3+}}$ is an oxidant ($\mathrm{Tl^{3+} + 2e^- \to Tl^+}$).
    \item \textbf{Group 14:} Stability of $+2$ state increases: $\mathrm{C(+4)} > \mathrm{Si(+4)} > \mathrm{Ge(+4)} > \mathrm{Sn(+4)} > \mathrm{Pb(+4)}$, whereas $\mathbf{\mathrm{Pb^{2+}}} > \mathbf{\mathrm{Sn^{2+}}} > \mathbf{\mathrm{Ge^{2+}}}$. Thus, $\mathrm{Sn^{2+}}$ is a reducing agent, while $\mathrm{PbO_2}$ is a strong oxidant.
    \item \textbf{Group 15:} Stability of $+3$ state increases: $\mathbf{\mathrm{Bi^{3+}}} > \mathbf{\mathrm{Sb^{3+}}} > \mathbf{\mathrm{As^{3+}}}$. $\mathrm{Bi_2O_5}$ and $\mathrm{BiF_5}$ are powerful oxidants because $\mathrm{Bi(+5)}$ readily reduces to $+3$.
  \end{itemize}
\end{itemize}

% =========================================================================
% SECTION 3: CHEMICAL BONDING & MOLECULAR STRUCTURE
% =========================================================================
\section{Chemical Bonding \& Molecular Structure (Unit 3)}

\subsection{Ionic Bonding, Lattice Energy \& Fajan's Rules}
Ionic bonding arises from omnidirectional electrostatic attraction between oppositely charged ions:
\begin{equation}
  E_{\text{coulomb}} = -\frac{1}{4\pi\varepsilon_0}\frac{z_+ z_- e^2}{r_0}
\end{equation}
Lattice energy ($\Delta_{\text{lattice}} H^\circ$) is the energy required to completely separate 1 mole of solid ionic crystal into its constituent gaseous ions at infinite separation.

\begin{conceptbox}[Fajan's Rules: Covalent Character in Ionic Bonds]
No chemical bond is $100\%$ ionic. When a cation approaches an anion, the cation pulls the valence electron cloud of the anion toward itself while repelling its positive nucleus, polarizing the anion. High polarization introduces \textbf{covalent character}.
\begin{enumerate}
  \item \textbf{Small Cation Size:} Smaller cations possess higher positive charge density and greater polarizing power ($\phi = \text{charge}/\text{radius}$):
  \begin{equation}
    \text{Covalent character: } \mathrm{LiCl > NaCl > KCl > RbCl > CsCl} \implies \text{Melting Point: } \mathrm{NaCl > LiCl}
  \end{equation}
  \item \textbf{Large Anion Size:} Larger anions have diffuse electron clouds that are easily distorted (high polarizability):
  \begin{equation}
    \text{Covalent character: } \mathrm{LiI > LiBr > LiCl > LiF}
  \end{equation}
  \item \textbf{High Ionic Charges:} High positive charge on cation and high negative charge on anion dramatically maximize polarization:
  \begin{equation}
    \text{Covalent character: } \mathrm{SnCl_4 (Sn^{4+}) > SnCl_2 (Sn^{2+})}, \qquad \mathrm{AlCl_3 (Al^{3+}) > MgCl_2 (Mg^{2+}) > NaCl (Na^+)}
  \end{equation}
  \item \textbf{Pseudo-Inert Gas Electron Configuration ($ns^2 np^6 nd^{10}$):}
  A cation with 18 valence electrons ($\mathrm{Cu^+, Ag^+, Zn^{2+}, Cd^{2+}}$) exerts far greater polarizing power than a noble-gas configuration cation ($ns^2 np^6$, 8 electrons; e.g., $\mathrm{Na^+, K^+, Ca^{2+}}$) of comparable size and charge, because $d$-electrons shield the nuclear charge very poorly:
  \begin{equation}
    \mathrm{CuCl\ (18e^-)} \text{ is far more covalent and insoluble than } \mathrm{NaCl\ (8e^-)} \quad [r(\mathrm{Cu^+}) \approx r(\mathrm{Na^+})]
  \end{equation}
\end{enumerate}
\end{conceptbox}

\subsection{Lewis Structures \& Formal Charge}
\begin{formulabox}[Formal Charge Calculation]
Formal charge represents the hypothetical electric charge assigned to an individual atom in a Lewis structure assuming equal sharing of bonding electrons:
\begin{equation}
  \text{FC} = V - L - \frac{1}{2}S
\end{equation}
where $V$ is the number of valence electrons of the isolated neutral atom, $L$ is the number of non-bonding lone-pair electrons, and $S$ is the number of shared bonding electrons.
\textbf{Example: Ozone ($\mathrm{O_3}$):}
\begin{itemize}
  \item Central Oxygen atom ($1$ double bond, $1$ single coordinate bond, $1$ lone pair):
  \begin{equation}
    \text{FC} = 6 - 2 - \frac{1}{2}(6) = +1
  \end{equation}
  \item Double-bonded terminal Oxygen ($2$ lone pairs, $4$ shared electrons):
  \begin{equation}
    \text{FC} = 6 - 4 - \frac{1}{2}(4) = 0
  \end{equation}
  \item Single-bonded terminal Oxygen ($3$ lone pairs, $2$ shared electrons):
  \begin{equation}
    \text{FC} = 6 - 6 - \frac{1}{2}(2) = -1
  \end{equation}
\end{itemize}
\end{formulabox}

\subsection{Dipole Moments \& Vector Summation}
The electric dipole moment $\vec{\mu}$ measures bond and molecular polarity:
\begin{equation}
  \mu = q \times d \quad [1\,\text{Debye (D)} = 3.33564 \times 10^{-30}\,\si{C.m}]
\end{equation}
Molecular dipole moment is the vector sum of all individual bond moments and lone-pair moments:
\begin{equation}
  \mu_{\text{net}} = \sqrt{\mu_1^2 + \mu_2^2 + 2\mu_1\mu_2\cos\theta}
\end{equation}

\begin{trapbox}[The $\mathrm{NH_3}$ vs $\mathrm{NF_3}$ Dipole Moment Paradox]
Both Ammonia ($\mathrm{NH_3}$) and Nitrogen Trifluoride ($\mathrm{NF_3}$) exhibit identical trigonal pyramidal geometry ($SN = 4$, 1 lone pair). However:
\begin{equation}
  \mu(\mathrm{NH_3}) = 1.47\,\si{D} \gg \mu(\mathrm{NF_3}) = 0.23\,\si{D}
\end{equation}
\textbf{Rigorous Vector Explanation:}
\begin{itemize}
  \item In $\mathrm{NH_3}$, Nitrogen ($\chi = 3.0$) is more electronegative than Hydrogen ($\chi = 2.1$). The three individual $\mathrm{N-H}$ bond dipoles point upward toward Nitrogen, \textbf{reinforcing and adding constructively} to the upward orbital dipole of the localized lone pair.
  \item In $\mathrm{NF_3}$, Fluorine ($\chi = 4.0$) is far more electronegative than Nitrogen ($\chi = 3.0$). The three individual $\mathrm{N-F}$ bond dipoles point downward away from Nitrogen, \textbf{opposing and destructively cancelling} the upward lone-pair dipole.
\end{itemize}
\end{trapbox}

\subsection{VSEPR Theory: Master 3D Geometry Framework}
Valence Shell Electron Pair Repulsion (VSEPR) theory dictates that electron pairs surrounding a central atom repel one another and adopt spatial orientations that minimize electrostatic repulsion.
\begin{equation}
  \text{Repulsion Hierarchy: } lp-lp > lp-bp > bp-bp
\end{equation}
Steric Number ($SN$) is computed by:
\begin{equation}
  SN = (\text{Number of }\sigma\text{-bonds}) + (\text{Number of localized lone pairs on central atom})
\end{equation}

\begin{table}[htbp]
\centering
\small
\renewcommand{\arraystretch}{1.25}
\begin{tabularx}{\textwidth}{c c c p{3.2cm} p{2.8cm} Y}
\toprule
$SN$ & $\sigma\text{-bonds}$ & Lone Pairs & \textbf{Electron Geometry} & \textbf{Molecular Geometry} & \textbf{High-Yield Archetypes \& Angles} \\
\midrule
$2$ & $2$ & $0$ & Linear & Linear & $\mathrm{BeCl_2, CO_2, HCN, HgCl_2}$ ($180^\circ$) \\
\midrule
$3$ & $3$ & $0$ & Trigonal Planar & Trigonal Planar & $\mathrm{BF_3, BCl_3, SO_3, NO_3^-}$ ($120^\circ$) \\
$3$ & $2$ & $1$ & Trigonal Planar & Bent / Angular & $\mathrm{SO_2, O_3, NO_2^-, SnCl_2}$ ($< 120^\circ$, $\approx 119^\circ$) \\
\midrule
$4$ & $4$ & $0$ & Tetrahedral & Tetrahedral & $\mathrm{CH_4, CCl_4, NH_4^+, SO_4^{2-}}$ ($109^\circ 28'$) \\
$4$ & $3$ & $1$ & Tetrahedral & Trigonal Pyramidal & $\mathrm{NH_3, PCl_3, H_3O^+, ClO_3^-}$ ($107^\circ$) \\
$4$ & $2$ & $2$ & Tetrahedral & Bent / V-shaped & $\mathrm{H_2O, H_2S, OF_2, SCl_2}$ ($104.5^\circ$) \\
\midrule
$5$ & $5$ & $0$ & Trigonal Bipyramidal & Trigonal Bipyramidal & $\mathrm{PCl_5, PF_5, AsF_5}$ (eq $120^\circ$, ax $90^\circ$) \\
$5$ & $4$ & $1$ & Trigonal Bipyramidal & Seesaw & $\mathrm{SF_4, SeF_4, TeCl_4}$ ($173^\circ, 102^\circ$) \\
$5$ & $3$ & $2$ & Trigonal Bipyramidal & T-shaped & $\mathrm{ClF_3, BrF_3, IF_3}$ ($87.5^\circ$) \\
$5$ & $2$ & $3$ & Trigonal Bipyramidal & Linear & $\mathrm{XeF_2, I_3^-, ICl_2^-}$ ($180^\circ$) \\
\midrule
$6$ & $6$ & $0$ & Octahedral & Octahedral & $\mathrm{SF_6, [PF_6]^-, [AlF_6]^{3-}}$ ($90^\circ$) \\
$6$ & $5$ & $1$ & Octahedral & Square Pyramidal & $\mathrm{BrF_5, IF_5, XeOF_4}$ ($86^\circ$) \\
$6$ & $4$ & $2$ & Octahedral & Square Planar & $\mathrm{XeF_4, [ICl_4]^-, [BrF_4]^-}$ ($90^\circ$) \\
\midrule
$7$ & $7$ & $0$ & Pentagonal Bipyramidal & Pentagonal Bipyramidal & $\mathrm{IF_7}$ (eq $72^\circ$, ax $90^\circ$) \\
$7$ & $6$ & $1$ & Pentagonal Bipyramidal & Distorted Octahedral & $\mathrm{XeF_6}$ (fluxional lone pair in face) \\
\bottomrule
\end{tabularx}
\caption{Master VSEPR Molecular Geometry Classification Table ($SN = 2$ to $7$).}
\end{table}

\begin{conceptbox}[Bent's Rule \& Bond Length Asymmetry in $\mathrm{PCl_5}$]
\begin{itemize}
  \item \textbf{Bent's Rule Formulation:} More electronegative substituents prefer hybrid orbitals with \textbf{less $s$-character} (more $p$-character). Lone pairs and electropositive substituents prefer hybrid orbitals with \textbf{more $s$-character}.
  \item \textbf{Application to Trigonal Bipyramidal Geometry ($sp^3d = sp^2 + pd$):}
  \begin{itemize}
    \item Equatorial hybrid orbitals ($sp^2$) possess $33.3\%$ $s$-character.
    \item Axial hybrid orbitals ($pd$) possess $0\%$ $s$-character ($50\% p + 50\% d_{z^2}$).
  \end{itemize}
  \item \textbf{Asymmetry in Gaseous $\mathrm{PCl_5}$:}
  \begin{itemize}
    \item The two axial $\mathrm{P-Cl}$ bonds experience intense $90^\circ$ repulsion from three equatorial bonds.
    \item Consequently, \textbf{axial bonds are longer and weaker} ($214\,\si{pm}$) than equatorial bonds ($202\,\si{pm}$).
    \item Upon gentle heating, $\mathrm{PCl_5}$ easily dissociates its two axial chlorines: $\mathrm{PCl_5(g) \xrightarrow{\Delta} PCl_3(g) + Cl_2(g)}$.
    \item In the crystalline solid phase, $\mathrm{PCl_5}$ ionizes into ionic tetrahedral and octahedral entities: $\mathbf{2\mathrm{PCl_5(s)} \longrightarrow [\mathrm{PCl_4}]^+[\mathrm{PCl_6}]^-}$. Solid $\mathrm{PBr_5}$ exists as $[\mathrm{PBr_4}]^+\mathrm{Br^-}$.
  \end{itemize}
  \item \textbf{Drago's Rule:}
  In hydrides of Group 15 and 16 elements from the 3rd period onwards ($\mathrm{PH_3, AsH_3, SbH_3, H_2S, H_2Se, H_2Te}$), the central atom does \textbf{NOT} undergo hybridization. Bonding occurs via nearly pure $p$-orbitals at $\approx 90^\circ$ angles ($\mathrm{PH_3}: 93.5^\circ, \mathrm{H_2S}: 92^\circ$). The lone pair resides in an inert, spherically symmetric $s$-orbital with negligible Lewis basicity.
\end{itemize}
\end{conceptbox}

\begin{usecasebox}[High-Yield Problem Archetype: Systematic VSEPR Steric Number Prediction]
\textbf{Predicting Hybridization and 3D Geometry for Problematic NEET Species:}
Formula for Steric Number of central atom $A$ in $AB_n^{q}$:
\begin{equation*}
  SN = \frac{1}{2}\left[V + M - C + A\right]
\end{equation*}
where $V = \text{valence } e^- \text{ of } A$, $M = \text{monovalent atoms attached } (\mathrm{H, F, Cl, Br, I})$, $C = \text{cationic charge}$, $A = \text{anionic charge}$.
\begin{enumerate}
  \item \textbf{Triiodide Ion ($\mathrm{I_3^-}$):} Central $\mathrm{I}$ has $V = 7$, two attached monovalent iodines ($M = 2$), anionic charge $A = 1$:
  \begin{equation*}
    SN = \frac{1}{2}[7 + 2 - 0 + 1] = 5 \implies sp^3d \text{ hybridization}
  \end{equation*}
  With 2 bond pairs and $5 - 2 = 3$ lone pairs: all three lone pairs occupy equatorial positions at $120^\circ$ $\implies$ \textbf{Linear Geometry} (bond angle $180^\circ$).
  \item \textbf{Xenon Tetrafluoride ($\mathrm{XeF_4}$):} Central $\mathrm{Xe}$ has $V = 8$, $M = 4$:
  \begin{equation*}
    SN = \frac{1}{2}[8 + 4] = 6 \implies sp^3d^2 \text{ hybridization}
  \end{equation*}
  With 4 bond pairs and $6 - 4 = 2$ lone pairs: lone pairs occupy opposite trans-axial positions at $180^\circ$ $\implies$ \textbf{Square Planar Geometry} ($90^\circ$).
  \item \textbf{Sulfur Tetrafluoride ($\mathrm{SF_4}$):} Central $\mathrm{S}$ has $V = 6$, $M = 4$:
  \begin{equation*}
    SN = \frac{1}{2}[6 + 4] = 5 \implies sp^3d \text{ hybridization}
  \end{equation*}
  With 4 bond pairs and 1 lone pair: lone pair occupies an equatorial position $\implies$ \textbf{Seesaw Geometry}.
  \item \textbf{Chlorine Trifluoride ($\mathrm{ClF_3}$):} Central $\mathrm{Cl}$ has $V = 7$, $M = 3$:
  \begin{equation*}
    SN = \frac{1}{2}[7 + 3] = 5 \implies sp^3d \text{ hybridization}
  \end{equation*}
  With 3 bond pairs and 2 lone pairs: both lone pairs occupy equatorial positions $\implies$ \textbf{T-shaped Geometry}.
\end{enumerate}
\end{usecasebox}

\subsection{Molecular Orbital Theory (MOT): Energy Levels, Bond Orders \& Magnetism}
Molecular Orbital Theory constructs molecular orbitals by Linear Combination of Atomic Orbitals (LCAO):
\begin{equation}
  \psi_{\text{MO}} = c_A \psi_A \pm c_B \psi_B \implies \psi_b = \frac{1}{\sqrt{2}}(\psi_A + \psi_B), \qquad \psi_a = \frac{1}{\sqrt{2}}(\psi_A - \psi_B)
\end{equation}

\subsubsection{1. Diatomic Molecules with $\le 14$ Valence Electrons (with $sp$-Mixing: $\mathrm{B_2, C_2, N_2}$)}
Because the energy separation between $2s$ and $2p_z$ atomic orbitals is small in $\mathrm{B, C, N}$, quantum mechanical mixing occurs between $\sigma 2s$ and $\sigma 2p_z$. This raises the energy of $\sigma 2p_z$ above the degenerate $\pi 2p_x, \pi 2p_y$ orbitals:
\begin{equation}
  \sigma 1s < \sigma^* 1s < \sigma 2s < \sigma^* 2s < (\pi 2p_x = \pi 2p_y) < \sigma 2p_z < (\pi^* 2p_x = \pi^* 2p_y) < \sigma^* 2p_z
\end{equation}

\subsubsection{2. Diatomic Molecules with $> 14$ Valence Electrons (without $sp$-Mixing: $\mathrm{O_2, F_2}$)}
In Oxygen and Fluorine, higher nuclear charge depresses the $2s$ level, creating a large $\Delta E(2s-2p_z)$ gap ($> 15\,\si{eV}$) that precludes $sp$-mixing. Normal energy ordering applies:
\begin{equation}
  \sigma 1s < \sigma^* 1s < \sigma 2s < \sigma^* 2s < \sigma 2p_z < (\pi 2p_x = \pi 2p_y) < (\pi^* 2p_x = \pi^* 2p_y) < \sigma^* 2p_z
\end{equation}

\begin{formulabox}[Bond Order \& Magnetic Properties Formulation]
\begin{equation}
  \text{Bond Order (BO)} = \frac{N_b - N_a}{2}
\end{equation}
where $N_b$ is the number of electrons in bonding MOs, and $N_a$ is the number of electrons in antibonding MOs.
\begin{itemize}
  \item If all electrons are paired $\implies$ \textbf{Diamagnetic}.
  \item If one or more unpaired electrons exist $\implies$ \textbf{Paramagnetic} ($\mu_s = \sqrt{n(n+2)}\,\si{BM}$).
\end{itemize}
\end{formulabox}

\begin{table}[htbp]
\centering
\footnotesize
\renewcommand{\arraystretch}{1.25}
\begin{tabularx}{\textwidth}{c c p{6.8cm} c c L}
\toprule
\textbf{Species} & \textbf{Total $e^-$} & \textbf{Valence MO Configuration} & $N_b / N_a$ & \textbf{BO} & \textbf{Magnetic State \& Properties} \\
\midrule
$\mathrm{B_2}$ & $10$ & $(\sigma 2s)^2 (\sigma^* 2s)^2 (\pi 2p_x)^1 (\pi 2p_y)^1$ & $6 / 4$ & $1.0$ & \textbf{Paramagnetic} ($2$ unpaired $e^-$ in $\pi$) \\
$\mathrm{C_2}$ & $12$ & $(\sigma 2s)^2 (\sigma^* 2s)^2 (\pi 2p_x)^2 (\pi 2p_y)^2$ & $8 / 4$ & $2.0$ & \textbf{Diamagnetic} (\textbf{Both bonds are $\pi$-bonds!}) \\
$\mathrm{N_2}$ & $14$ & $(\sigma 2s)^2 (\sigma^* 2s)^2 (\pi 2p_x)^2 (\pi 2p_y)^2 (\sigma 2p_z)^2$ & $10 / 4$ & $3.0$ & \textbf{Diamagnetic} (Triple bond: $1\sigma + 2\pi$, max stability) \\
$\mathrm{N_2^+}$ & $13$ & $(\sigma 2s)^2 (\sigma^* 2s)^2 (\pi 2p_x)^2 (\pi 2p_y)^2 (\sigma 2p_z)^1$ & $9 / 4$ & $2.5$ & \textbf{Paramagnetic} ($1$ unpaired $e^-$ in $\sigma 2p_z$) \\
$\mathrm{O_2^+}$ & $15$ & $(\sigma 2s)^2 (\sigma^* 2s)^2 (\sigma 2p_z)^2 (\pi 2p)^4 (\pi^* 2p_x)^1$ & $10 / 5$ & $2.5$ & \textbf{Paramagnetic} (Highest BO among oxygen species) \\
$\mathrm{O_2}$ & $16$ & $(\sigma 2s)^2 (\sigma^* 2s)^2 (\sigma 2p_z)^2 (\pi 2p)^4 (\pi^* 2p_x)^1 (\pi^* 2p_y)^1$ & $10 / 6$ & $2.0$ & \textbf{Paramagnetic} ($2$ unpaired $e^-$ in degenerate $\pi^*$) \\
$\mathrm{O_2^-}$ & $17$ & $(\sigma 2s)^2 (\sigma^* 2s)^2 (\sigma 2p_z)^2 (\pi 2p)^4 (\pi^* 2p_x)^2 (\pi^* 2p_y)^1$ & $10 / 7$ & $1.5$ & \textbf{Paramagnetic} (Superoxide ion) \\
$\mathrm{O_2^{2-}}$ & $18$ & $(\sigma 2s)^2 (\sigma^* 2s)^2 (\sigma 2p_z)^2 (\pi 2p)^4 (\pi^* 2p_x)^2 (\pi^* 2p_y)^2$ & $10 / 8$ & $1.0$ & \textbf{Diamagnetic} (Peroxide ion, single bond) \\
\bottomrule
\end{tabularx}
\caption{Comprehensive Molecular Orbital Configurations, Bond Orders, and Magnetism.}
\end{table}

\begin{ncertbox}[NCERT High-Yield Benchmark: Oxygen Species Comparison]
\begin{itemize}
  \item \textbf{Bond Order Sequence:} $\mathrm{O_2^+ (2.5) > O_2 (2.0) > O_2^- (1.5) > O_2^{2-} (1.0)}$.
  \item \textbf{Bond Length Sequence (Inversely proportional to Bond Order):}
  \begin{equation}
    \mathbf{\mathrm{O_2^+ (112\,pm) < O_2 (121\,pm) < O_2^- (128\,pm) < O_2^{2-} (149\,pm)}}
  \end{equation}
  \item \textbf{Bond Dissociation Energy Sequence:} $\mathrm{O_2^+ > O_2 > O_2^- > O_2^{2-}}$.
\end{itemize}
\end{ncertbox}

\subsection{Hydrogen Bonding: Intermolecular vs Intramolecular}
Hydrogen bonding is an electrostatic dipole-dipole attraction occurring between a covalently bound Hydrogen atom (attached to highly electronegative $\mathrm{F, O, N}$) and another electronegative atom with a localized lone pair.
\begin{itemize}
  \item \textbf{Intermolecular Hydrogen Bonding:} Occurs between distinct molecules (e.g., $\mathrm{H_2O}$, $\mathrm{HF}$, $\mathrm{NH_3}$, $\mathrm{C_2H_5OH}$, \textit{p}-nitrophenol).
  \textbf{Effects:} Elevated boiling points, higher viscosity, lower volatility, and high water solubility. Ice has an open cage-like tetrahedral structure making it less dense than liquid water at $4^\circ\text{C}$.
  \item \textbf{Intramolecular Hydrogen Bonding (Chelation):} Occurs within the same molecule, forming a 5- or 6-membered ring (e.g., \textit{o}-nitrophenol, salicylaldehyde, salicylic acid).
  \textbf{Effects:} Prevents intermolecular association; results in \textbf{significantly lower boiling points}, higher volatility, and \textbf{steam volatility} (enabling clean separation of \textit{o}-nitrophenol from \textit{p}-nitrophenol via steam distillation).
\end{itemize}

% =========================================================================
% SECTION 4: COORDINATION CHEMISTRY
% =========================================================================
\section{Coordination Chemistry (Unit 12)}

\subsection{Werner's Coordination Theory}
Alfred Werner (1893) postulated that transition metals exhibit two distinct types of valency:
\begin{enumerate}
  \item \textbf{Primary Valency (Ionizable):} Corresponds to the oxidation state of the metal ion. It is non-directional and satisfied strictly by negative anions.
  \item \textbf{Secondary Valency (Non-Ionizable):} Corresponds to the coordination number of the metal ion ($4$ or $6$). It is strictly directional in 3D space (establishing spatial geometry) and satisfied by neutral molecules or negative ligands.
\end{enumerate}

\begin{table}[htbp]
\centering
\small
\renewcommand{\arraystretch}{1.25}
\begin{tabularx}{\textwidth}{p{3.2cm} p{3.2cm} c c c}
\toprule
\textbf{Cobalt Ammine Complex} & \textbf{Modern IUPAC Formulation} & \textbf{Moles $\mathrm{AgCl\downarrow}$} & \textbf{Total Ions} & \textbf{Electrical Conductivity} \\
\midrule
$\mathrm{CoCl_3\cdot 6NH_3}$ & $[\mathrm{Co(NH_3)_6}]Cl_3$ & $3$ & $4$ ($1:3$ electrolyte) & Highest ($430\,\si{S.cm^2.mol^{-1}}$) \\
$\mathrm{CoCl_3\cdot 5NH_3}$ & $[\mathrm{Co(NH_3)_5Cl}]Cl_2$ & $2$ & $3$ ($1:2$ electrolyte) & Moderate ($260\,\si{S.cm^2.mol^{-1}}$) \\
$\mathrm{CoCl_3\cdot 4NH_3}$ & $[\mathrm{Co(NH_3)_4Cl_2}]Cl$ & $1$ & $2$ ($1:1$ electrolyte) & Low ($100\,\si{S.cm^2.mol^{-1}}$) \\
$\mathrm{CoCl_3\cdot 3NH_3}$ & $[\mathrm{Co(NH_3)_3Cl_3}]$ & $0$ & $0$ (Neutral molecule) & Non-conducting ($0$) \\
\bottomrule
\end{tabularx}
\caption{Werner Formulation, Precipitation Stoichiometry, and Solution Conductivity.}
\end{table}

\subsection{Nomenclature \& Isomerism in Coordination Entities}
\subsubsection{IUPAC Nomenclature Principles}
\begin{enumerate}
  \item Cation is named first in both positive and negative coordination entities.
  \item Ligands are listed in strict alphabetical order before the central metal name, regardless of charge.
  \item Multiplicative prefixes: \textit{di-, tri-, tetra-} for simple ligands; \textit{bis-, tris-, tetrakis-} for polydentate or ligands containing prefixes (e.g., \textit{ethylenediamine}).
  \item Anionic ligands end in \textit{-ido} ($\mathrm{Cl^-} \to \text{chlorido}, \mathrm{CN^-} \to \text{cyanido}, \mathrm{OH^-} \to \text{hydroxido}$).
  \item If the complex sphere is anionic, the metal name ends in \textbf{\textit{-ate}} ($\mathrm{Fe} \to \text{ferrate}, \mathrm{Cu} \to \text{cuprate}, \mathrm{Ag} \to \text{argentate}$).
  \item Oxidation state of the metal is designated by Roman numerals in parentheses: e.g., $\mathrm{K_4[Fe(CN)_6]} \implies \text{Potassium hexacyanidoferrate(II)}$.
\end{enumerate}

\subsubsection{Structural Isomerism}
\begin{itemize}
  \item \textbf{Ionization Isomerism:} Exchange of ionizable counter-ion with coordinated ligand:
  \begin{equation}
    [\mathrm{Co(NH_3)_5SO_4}]\mathrm{Br}\text{ (gives }\mathrm{AgBr\downarrow}\text{ with }\mathrm{AgNO_3}) \quad \text{vs} \quad [\mathrm{Co(NH_3)_5Br}]\mathrm{SO_4}\text{ (gives }\mathrm{BaSO_4\downarrow}\text{ with }\mathrm{BaCl_2})
  \end{equation}
  \item \textbf{Solvate / Hydrate Isomerism:} Water acts either as a coordinated ligand or crystal lattice water:
  \begin{equation}
    [\mathrm{Cr(H_2O)_6}]Cl_3\text{ (Violet, 3 }\mathrm{Cl^-}) \quad \text{vs} \quad [\mathrm{Cr(H_2O)_5Cl}]Cl_2\cdot \mathrm{H_2O}\text{ (Grey-green, 2 }\mathrm{Cl^-})
  \end{equation}
  \item \textbf{Linkage Isomerism:} Arises with ambidentate ligands having two distinct coordinating donor atoms:
  \begin{equation}
    [\mathrm{Co(NH_3)_5(NO_2)}]^{2+}\text{ (Nitro, yellow-brown, via N)} \quad \text{vs} \quad [\mathrm{Co(NH_3)_5(ONO)}]^{2+}\text{ (Nitrito, red, via O)}
  \end{equation}
  \item \textbf{Coordination Isomerism:} Interchange of ligands between cationic and anionic complex entities:
  \begin{equation}
    [\mathrm{Co(NH_3)_6}][\mathrm{Cr(CN)_6}] \quad \text{vs} \quad [\mathrm{Cr(NH_3)_6}][\mathrm{Co(CN)_6}]
  \end{equation}
\end{itemize}

\subsubsection{Stereoisomerism}
\begin{enumerate}
  \item \textbf{Geometrical Isomerism (GI):}
  \begin{itemize}
    \item \textbf{Square Planar ($CN=4$):} $[MA_2B_2]$ exhibits \textit{cis} and \textit{trans} (e.g., \textit{cis}-platin $[\mathrm{Pt(NH_3)_2Cl_2}]$ antitumor drug vs inactive \textit{trans}). $[MABCD]$ yields 3 distinct geometrical isomers. Tetrahedral complexes \textbf{NEVER} show GI because all 4 positions are mutually adjacent ($109.5^\circ$).
    \item \textbf{Octahedral ($CN=6$):}
    \begin{itemize}
      \item $[MA_4B_2]$: \textit{cis} ($90^\circ$) and \textit{trans} ($180^\circ$).
      \item $[MA_3B_3]$: \textbf{Facial (\textit{fac})} where three identical ligands occupy corners of the same octahedral triangular face, and \textbf{Meridional (\textit{mer})} where ligands occupy a meridian plane surrounding the metal.
    \end{itemize}
  \end{itemize}
  \item \textbf{Optical Isomerism (Chirality):}
  Molecules lacking an improper axis of rotation ($S_n$, including plane of symmetry $\sigma$ and inversion centre $i$) exist as non-superimposable mirror image enantiomers ($d$ and $l$).
  \begin{itemize}
    \item $[M(AA)_3]$ (e.g., $[Co(en)_3]^{3+}$) is chiral and resolves into $\Delta$ and $\Lambda$ optical enantiomers.
    \item $[M(AA)_2B_2]$ (e.g., $[Co(en)_2Cl_2]^+$):
    \begin{itemize}
      \item \textbf{\textit{cis}-isomer:} Asymmetric, lacks plane of symmetry $\implies$ \textbf{Optically Active} ($2$ enantiomers).
      \item \textbf{\textit{trans}-isomer:} Possesses a plane of symmetry ($\sigma_h$) $\implies$ \textbf{Optically Inactive}.
    \end{itemize}
    \item Square planar complexes have a horizontal plane of symmetry and \textbf{NEVER} show optical isomerism.
  \end{itemize}
\end{enumerate}

\subsection{Valence Bond Theory (VBT) \& Magnetic States}
Linus Pauling applied hybridization of metal $d, s, p$ orbitals to explain coordinate bonding.
\begin{itemize}
  \item \textbf{Inner Orbital Complexes ($d^2sp^3$):} Strong-field ligands force electron pairing in $(n-1)d$ orbitals, providing two empty inner $d$-orbitals (low-spin / spin-paired, e.g., $[\mathrm{Fe(CN)_6}]^{4-}$: $3d^6$ paired, diamagnetic, $\mu = 0$).
  \item \textbf{Outer Orbital Complexes ($sp^3d^2$):} Weak-field ligands cannot force pairing; vacant outer $nd$ orbitals are hybridized (high-spin / spin-free, e.g., $[\mathrm{Fe(H_2O)_6}]^{2+}$: $3d^6$ with 4 unpaired electrons, paramagnetic, $\mu = \sqrt{4(6)} = 4.90\,\si{BM}$).
  \item \textbf{Limitations of VBT:} Fails to predict absorption spectra or color; cannot quantitatively explain temperature-dependent magnetic susceptibilities; cannot explain why identical $d^8$ configuration forms tetrahedral paramagnetic $[\mathrm{NiCl_4}]^{2-}$ but square planar diamagnetic $[\mathrm{Ni(CN)_4}]^{2-}$.
\end{itemize}

\subsection{Crystal Field Theory (CFT): Splitting Diagrams \& CFSE}
CFT treats the metal-ligand interaction as a purely electrostatic point-charge perturbation. In an isolated gaseous ion, all five $d$-orbitals are degenerate. Surrounding ligand fields split this degeneracy:

\subsubsection{1. Octahedral Crystal Field Splitting ($\Delta_o$)}
Six ligands approach along the Cartesian $\pm x, \pm y, \pm z$ axes:
\begin{itemize}
  \item The $e_g$ orbitals ($d_{x^2-y^2}, d_{z^2}$) point directly along the axes $\implies$ experience maximal electrostatic repulsion $\implies$ destabilized and raised by $\mathbf{+0.6\Delta_o}$ ($+\frac{3}{5}\Delta_o$).
  \item The $t_{2g}$ orbitals ($d_{xy}, d_{yz}, d_{zx}$) point between the Cartesian axes $\implies$ experience less repulsion $\implies$ stabilized and lowered by $\mathbf{-0.4\Delta_o}$ ($-\frac{2}{5}\Delta_o$).
  \item The barycentre (weighted average energy) remains invariant: $3(-0.4\Delta_o) + 2(+0.6\Delta_o) = 0$.
\end{itemize}

\subsubsection{2. Tetrahedral Crystal Field Splitting ($\Delta_t$)}
Four ligands approach the alternate corners of a cube containing the metal at the centre:
\begin{itemize}
  \item The $t_2$ orbitals ($d_{xy}, d_{yz}, d_{zx}$) lie closer to ligand directions $\implies$ raised by $\mathbf{+0.4\Delta_t}$.
  \item The $e$ orbitals ($d_{x^2-y^2}, d_{z^2}$) lie further away $\implies$ lowered by $\mathbf{-0.6\Delta_t}$.
  \item \textbf{Splitting Magnitude Relation:}
  \begin{equation}
    \Delta_t = \frac{4}{9}\Delta_o
  \end{equation}
  Because $\Delta_t$ is less than half of $\Delta_o$, the crystal field splitting is \textbf{always smaller than pairing energy} ($\Delta_t < P$). Consequently, \textbf{tetrahedral complexes are ALWAYS high-spin}.
\end{itemize}

\begin{formulabox}[Spectrochemical Series \& CFSE Calculations]
Ligands arranged in order of increasing crystal field splitting magnitude $\Delta_o$:
\begin{align}
  \mathrm{I^- < Br^- < S^{2-} < SCN^- < Cl^- < N_3^- < F^- < OH^- < C_2O_4^{2-}} \nonumber \\
  {} < \mathrm{H_2O < NCS^- < edta^{4-} < NH_3 < en < NO_2^- < CN^- < CO}
\end{align}
\textbf{High-Spin vs Low-Spin Criterion for $d^4, d^5, d^6, d^7$ in Octahedral Fields:}
\begin{itemize}
  \item \textbf{Weak Field Ligands ($\Delta_o < P$):} Electron promotes to $e_g$ before pairing in $t_{2g}$ $\implies$ \textbf{High Spin}.
  \item \textbf{Strong Field Ligands ($\Delta_o > P$):} Electrons pair completely in $t_{2g}$ before occupying $e_g$ $\implies$ \textbf{Low Spin}.
\end{itemize}
\textbf{Octahedral CFSE Formula:}
\begin{equation}
  \text{CFSE} = \left(-0.4\,n_{t_{2g}} + 0.6\,n_{e_g}\right)\Delta_o + m P
\end{equation}
where $m$ is the number of extra electron pairs formed relative to the free gaseous ion.
\end{formulabox}

\begin{table}[htbp]
\centering
\small
\renewcommand{\arraystretch}{1.25}
\begin{tabularx}{\textwidth}{c p{3.2cm} c c p{3.2cm} c c}
\toprule
& \multicolumn{3}{c}{\textbf{High Spin (Weak Field: $\Delta_o < P$)}} & \multicolumn{3}{c}{\textbf{Low Spin (Strong Field: $\Delta_o > P$)}} \\
\cmidrule(lr){2-4} \cmidrule(lr){5-7}
$d^n$ & \textbf{Config} & \textbf{Unpaired} & \textbf{CFSE} & \textbf{Config} & \textbf{Unpaired} & \textbf{CFSE} \\
\midrule
$d^1$ & $t_{2g}^1 e_g^0$ & $1$ & $-0.4\Delta_o$ & $t_{2g}^1 e_g^0$ & $1$ & $-0.4\Delta_o$ \\
$d^2$ & $t_{2g}^2 e_g^0$ & $2$ & $-0.8\Delta_o$ & $t_{2g}^2 e_g^0$ & $2$ & $-0.8\Delta_o$ \\
$d^3$ & $t_{2g}^3 e_g^0$ & $3$ & $-1.2\Delta_o$ & $t_{2g}^3 e_g^0$ & $3$ & $-1.2\Delta_o$ \\
$d^4$ & $t_{2g}^3 e_g^1$ & $4$ & $-0.6\Delta_o$ & $t_{2g}^4 e_g^0$ & $2$ & $-1.6\Delta_o + 1P$ \\
$d^5$ & $t_{2g}^3 e_g^2$ & $5$ & $\mathbf{0.0\Delta_o}$ & $t_{2g}^5 e_g^0$ & $1$ & $-2.0\Delta_o + 2P$ \\
$d^6$ & $t_{2g}^4 e_g^2$ & $4$ & $-0.4\Delta_o$ & $t_{2g}^6 e_g^0$ & $0$ & $\mathbf{-2.4\Delta_o + 2P}$ \\
$d^7$ & $t_{2g}^5 e_g^2$ & $3$ & $-0.8\Delta_o$ & $t_{2g}^6 e_g^1$ & $1$ & $-1.8\Delta_o + 1P$ \\
$d^8$ & $t_{2g}^6 e_g^2$ & $2$ & $-1.2\Delta_o$ & $t_{2g}^6 e_g^2$ & $2$ & $-1.2\Delta_o$ \\
$d^9$ & $t_{2g}^6 e_g^3$ & $1$ & $-0.6\Delta_o$ & $t_{2g}^6 e_g^3$ & $1$ & $-0.6\Delta_o$ \\
$d^{10}$ & $t_{2g}^6 e_g^4$ & $0$ & $0.0\Delta_o$ & $t_{2g}^6 e_g^4$ & $0$ & $0.0\Delta_o$ \\
\bottomrule
\end{tabularx}
\caption{Systematic Crystal Field Stabilization Energies (CFSE) for Octahedral Complexes.}
\end{table}

\begin{usecasebox}[High-Yield Problem Archetype: Octahedral CFSE \& Spin-Only Magnetic Moments]
\textbf{Archetype Problem:} Compare the electronic configuration, CFSE, and magnetic moment ($\mu_s$) for:
(a) $[\mathrm{CoF_6}]^{3-}$ vs $[\mathrm{Co(en)_3}]^{3+}$, and (b) $[\mathrm{Fe(CN)_6}]^{3-}$ vs $[\mathrm{Fe(H_2O)_6}]^{3+}$.
\begin{enumerate}
  \item \textbf{Cobalt(III) Complexes ($3d^6$ system):}
  \begin{itemize}
    \item \textbf{$[\mathrm{CoF_6}]^{3-}$:} $\mathrm{F^-}$ is a weak-field ligand ($\Delta_o < P$) $\implies$ \textbf{High Spin}.
    \begin{align*}
      \text{Configuration} &= t_{2g}^4 e_g^2, \quad \text{Unpaired } e^- = 4 \\
      \text{CFSE} &= (-0.4 \times 4 + 0.6 \times 2)\Delta_o = -0.4\Delta_o \\
      \mu_s &= \sqrt{4(4+2)} = \sqrt{24} \approx 4.90\,\si{BM} \quad (\text{Paramagnetic})
    \end{align*}
    \item \textbf{$[\mathrm{Co(en)_3}]^{3-}$:} \textit{en} is a strong-field chelating ligand ($\Delta_o > P$) $\implies$ \textbf{Low Spin}.
    \begin{align*}
      \text{Configuration} &= t_{2g}^6 e_g^0, \quad \text{Unpaired } e^- = 0 \\
      \text{CFSE} &= (-0.4 \times 6)\Delta_o + 2P = -2.4\Delta_o + 2P \\
      \mu_s &= 0\,\si{BM} \quad (\text{Diamagnetic})
    \end{align*}
  \end{itemize}
  \item \textbf{Iron(III) Complexes ($3d^5$ system):}
  \begin{itemize}
    \item \textbf{$[\mathrm{Fe(H_2O)_6}]^{3+}$:} $\mathrm{H_2O}$ is weak-field $\implies t_{2g}^3 e_g^2 \implies \text{CFSE} = 0$, $n=5 \implies \mu_s = 5.92\,\si{BM}$.
    \item \textbf{$[\mathrm{Fe(CN)_6}]^{3-}$:} $\mathrm{CN^-}$ is strong-field $\implies t_{2g}^5 e_g^0 \implies \text{CFSE} = -2.0\Delta_o + 2P$, $n=1 \implies \mu_s = 1.73\,\si{BM}$.
  \end{itemize}
\end{enumerate}
\end{usecasebox}

\subsection{Jahn-Teller Distortion Basics}
The Jahn-Teller theorem states: \textit{Any non-linear molecular system in a degenerate electronic state is unstable and will undergo geometrical distortion that lowers its symmetry and removes the electronic degeneracy.}
\begin{itemize}
  \item \textbf{Significant / Strong Jahn-Teller Distortion:} Occurs when the \textbf{$e_g$ subshell is asymmetrically occupied} because $e_g$ orbitals point directly along the ligand bond axes:
  \begin{equation}
    \text{High-spin } d^4\ (t_{2g}^3 e_g^1: \mathrm{Cr^{2+}, Mn^{3+}}), \quad \text{Low-spin } d^7\ (t_{2g}^6 e_g^1: \mathrm{Co^{2+}}), \quad \mathbf{d^9\ (t_{2g}^6 e_g^3: \mathrm{Cu^{2+}})}
  \end{equation}
  \item \textbf{Tetragonal Elongation ($z$-out distortion):}
  In $[\mathrm{Cu(H_2O)_6}]^{2+}$, putting the odd electron in $d_{z^2}$ shields the axial ligands, causing the two axial $\mathrm{Cu-O}$ bonds ($240\,\si{pm}$) to elongate significantly relative to the four equatorial bonds ($195\,\si{pm}$).
\end{itemize}

\subsection{Synergic Bonding in Metal Carbonyls}
\begin{conceptbox}[Synergic Bonding Paradigm in Homoleptic Metal Carbonyls]
In metal carbonyls (e.g., $\mathrm{Ni(CO)_4, Fe(CO)_5, Cr(CO)_6}$), the metal-carbon bond comprises a self-reinforcing dual interaction:
\begin{enumerate}
  \item \textbf{$\sigma$-Bond Formation ($M \leftarrow \mathrm{CO}$):} Carbon donates its non-bonding lone pair into an empty hybrid orbital of the transition metal.
  \item \textbf{$\pi$-Backdonation ($M \to \mathrm{CO}$):} Filled non-bonding $d$-orbitals of the metal overlap with empty $\pi^*$ antibonding molecular orbitals of Carbon Monoxide.
\end{enumerate}
\textbf{Consequences of Synergism:}
\begin{itemize}
  \item Strengthens the $M-\mathrm{C}$ bond (increases bond order, shortens bond length).
  \item Populates the $\pi^*$ orbital of $\mathrm{CO}$, \textbf{weakening the $\mathrm{C-O}$ bond} (lowering $\mathrm{C-O}$ bond order and reducing $\mathrm{C-O}$ vibrational stretching frequency $\nu_{\mathrm{CO}}$ in IR spectroscopy).
  \item Greater negative charge on metal increases backdonation $\implies$ weakest $\mathrm{C-O}$ bond:
  \begin{equation}
    \text{IR frequency: } \mathrm{[Mn(CO)_6]^+ > [Cr(CO)_6] > [V(CO)_6]^-}
  \end{equation}
\end{itemize}
\end{conceptbox}

% =========================================================================
% SECTION 5: COMPARATIVE CHEMISTRY OF P-BLOCK ELEMENTS
% =========================================================================
\section{Comparative Chemistry of $p$-Block Elements (Groups 13 to 18)}

\subsection{Group 13: Boron Family}
\begin{itemize}
  \item \textbf{Anomalous Boron:} Non-metal, high melting point ($2453\,\si{K}$), maximum covalency of 4 (absence of $d$-orbitals).
  \item \textbf{Diborane ($\mathrm{B_2H_6}$):}
  Electron-deficient dimer with 12 valence electrons.
  \textbf{Bonding Structure:} Contains four terminal two-centre-two-electron ($2c-2e$) $\mathrm{B-H}$ bonds ($119\,\si{pm}$), and two bridging three-centre-two-electron ($3c-2e$) $\mathrm{B-H-B}$ \textbf{banana bonds} ($133\,\si{pm}$).
  \textbf{Synthesis \& Reactions:}
  \begin{align}
    4\mathrm{BF_3} + 3\mathrm{LiAlH_4} &\xrightarrow{\text{ether}} 2\mathrm{B_2H_6} + 3\mathrm{LiF} + 3\mathrm{AlF_3} \\
    \mathrm{B_2H_6} + 3\mathrm{O_2} &\longrightarrow \mathrm{B_2O_3} + 3\mathrm{H_2O} \quad (\Delta H = -1976\,\si{kJ.mol^{-1}}, \text{green-edged flame}) \\
    3\mathrm{B_2H_6} + 6\mathrm{NH_3} &\xrightarrow{450\,\si{K}} 2\mathrm{B_3N_3H_6}\text{ (\textbf{Borazine / Inorganic Benzene})} + 12\mathrm{H_2}
  \end{align}
  \item \textbf{Orthoboric Acid ($\mathrm{H_3BO_3}$):} Layer structure with planar $\mathrm{BO_3^{3-}}$ units held by H-bonds. Acts as a \textbf{monobasic weak Lewis acid} (accepts $\mathrm{OH^-}$ from water rather than donating $\mathrm{H}^+$):
  \begin{equation}
    \mathrm{B(OH)_3} + 2\mathrm{H_2O} \rightleftharpoons [\mathrm{B(OH)_4}]^- + \mathrm{H_3O}^+ \quad (\text{p}K_a = 9.25)
  \end{equation}
  Cannot be directly titrated with $\mathrm{NaOH}$, but addition of \textit{cis}-polyhydroxy compounds (glycerol, mannitol) forms stable chelate complexes, releasing proton quantitatively.
  \item \textbf{Lewis Acidity of Boron Trihalides:}
  \begin{equation}
    \mathbf{\mathrm{BI_3 > BBr_3 > BCl_3 > BF_3}}
  \end{equation}
  \textbf{Reason:} In $\mathrm{BF_3}$, vacant $2p$ orbital of Boron overlaps efficiently with filled $2p$ orbital of Fluorine via $p\pi-p\pi$ backbonding, neutralizing electron deficiency. In $\mathrm{BCl_3}$ ($2p-3p$) and $\mathrm{BI_3}$ ($2p-5p$), overlap is ineffective due to large size discrepancy.
\end{itemize}

\subsection{Group 14: Carbon Family}
\begin{itemize}
  \item \textbf{Allotropes of Carbon:}
  \begin{enumerate}
    \item \textbf{Diamond:} $sp^3$ hybridized, tetrahedral 3D giant covalent network, $\mathrm{C-C}$ bond length $154\,\si{pm}$, electrical insulator, hardest natural substance.
    \item \textbf{Graphite:} $sp^2$ hybridized, planar hexagonal layers ($\mathrm{C-C} = 141.5\,\si{pm}$) separated by $340\,\si{pm}$ held by weak van der Waals forces; fourth $\pi$-electron is delocalized over the sheet $\implies$ \textbf{electrical conductor} and solid lubricant. \textbf{Thermodynamically most stable allotrope} ($\Delta_f H^\circ = 0$).
    \item \textbf{Fullerenes ($\mathrm{C_{60}}$ Buckminsterfullerene):} Football-shaped cage of 60 carbon atoms comprising 20 six-membered rings and 12 five-membered rings.
  \end{enumerate}
  \item \textbf{Silicones:} Organosilicon polymers with repeating $[R_2\mathrm{SiO}]_n$ units synthesized by hydrolysis of dialkyldichlorosilanes ($\mathrm{R_2SiCl_2} \xrightarrow{\mathrm{H_2O}} \mathrm{R_2Si(OH)_2} \xrightarrow{\text{polymerization}} \text{silicone}$). Highly water-repellent, thermally stable, electrical insulators.
  \item \textbf{Silicates:} Basic structural unit is the tetrahedral $[\mathrm{SiO_4}]^{4-}$ entity:
  \begin{itemize}
    \item \textbf{Orthosilicates:} Discrete isolated $[\mathrm{SiO_4}]^{4-}$ (e.g., Zircon $\mathrm{ZrSiO_4}$, Phenacite $\mathrm{Be_2SiO_4}$).
    \item \textbf{Pyrosilicates:} Share 1 oxygen atom: $[\mathrm{Si_2O_7}]^{6-}$ (e.g., Thortveitite $\mathrm{Sc_2Si_2O_7}$).
    \item \textbf{Cyclic Silicates:} Share 2 oxygens per tetrahedron: $[\mathrm{Si_3O_9}]^{6-}$ or $[\mathrm{Si_6O_{18}}]^{12-}$ (e.g., Beryl $\mathrm{Be_3Al_2Si_6O_{18}}$).
    \item \textbf{Chain Silicates (Pyroxenes):} Share 2 oxygens forming linear chains: $[\mathrm{SiO_3}]_n^{2n-}$.
    \item \textbf{Sheet Silicates (Phyllosilicates):} Share 3 oxygens forming 2D sheets: $[\mathrm{Si_2O_5}]_n^{2n-}$ (e.g., Talc, Mica).
    \item \textbf{3D Framework Silicates (Tectosilicates):} Share all 4 oxygens: $\mathrm{SiO_2}$ (Quartz, Feldspar, Zeolites).
  \end{itemize}
  \item \textbf{Zeolites:} Microporous aluminosilicates with open 3D honeycomb channel structures. \textbf{ZSM-5} acts as a shape-selective catalyst converting alcohols directly into gasoline.
\end{itemize}

\subsection{Group 15: Nitrogen Family}
\begin{itemize}
  \item \textbf{Dinitrogen ($\mathrm{N_2}$):} Chemically inert at ambient conditions due to extraordinarily high bond dissociation energy of the $\mathrm{N\equiv N}$ triple bond ($941.4\,\si{kJ.mol^{-1}}$).
  \item \textbf{Hydrides of Group 15 ($\mathrm{NH_3, PH_3, AsH_3, SbH_3, BiH_3}$):}
  \begin{itemize}
    \item \textbf{Boiling Point:} $\mathrm{BiH_3 > SbH_3 > \mathbf{NH_3} > AsH_3 > PH_3}$ ($\mathrm{NH_3}$ abnormally high due to H-bonding; $\mathrm{PH_3}$ is lowest).
    \item \textbf{Basicity:} $\mathrm{NH_3 > PH_3 > AsH_3 > SbH_3 > BiH_3}$ (nitrogen has smallest size and highest lone-pair charge density).
    \item \textbf{Thermal Stability:} $\mathrm{NH_3 > PH_3 > AsH_3 > SbH_3 > BiH_3}$ ($\mathrm{M-H}$ bond length increases).
    \item \textbf{Reducing Power:} $\mathrm{BiH_3 > SbH_3 > AsH_3 > PH_3 > NH_3}$ ($\mathrm{BiH_3}$ is the strongest reducing agent).
    \item \textbf{Bond Angle:} $\mathrm{NH_3 (107.8^\circ) > PH_3 (93.5^\circ) > AsH_3 (91.8^\circ) > SbH_3 (91.3^\circ)}$ (Drago's rule).
  \end{itemize}
  \item \textbf{Nitric Acid ($\mathrm{HNO_3}$ - Ostwald Process):}
  \begin{equation}
    4\mathrm{NH_3} + 5\mathrm{O_2} \xrightarrow[\mathrm{Pt/Rh}, 500\,\si{K}]{9\,\si{bar}} 4\mathrm{NO} + 6\mathrm{H_2O}, \quad 2\mathrm{NO} + \mathrm{O_2} \to 2\mathrm{NO_2}, \quad 3\mathrm{NO_2} + \mathrm{H_2O} \to 2\mathrm{HNO_3} + \mathrm{NO}
  \end{equation}
  \textbf{Reaction with Metals (NEET Benchmark):}
  \begin{align}
    \text{Copper with Dilute } \mathrm{HNO_3} &: 3\mathrm{Cu} + 8\mathrm{HNO_3(dil)} \longrightarrow 3\mathrm{Cu(NO_3)_2} + \mathbf{2\mathrm{NO\uparrow}} + 4\mathrm{H_2O} \\
    \text{Copper with Conc. } \mathrm{HNO_3} &: \mathrm{Cu} + 4\mathrm{HNO_3(conc)} \longrightarrow \mathrm{Cu(NO_3)_2} + \mathbf{2\mathrm{NO_2\uparrow}} + 2\mathrm{H_2O} \\
    \text{Zinc with Dilute } \mathrm{HNO_3} &: 4\mathrm{Zn} + 10\mathrm{HNO_3(dil)} \longrightarrow 4\mathrm{Zn(NO_3)_2} + \mathbf{\mathrm{N_2O\uparrow}} + 5\mathrm{H_2O} \\
    \text{Zinc with Conc. } \mathrm{HNO_3} &: \mathrm{Zn} + 4\mathrm{HNO_3(conc)} \longrightarrow \mathrm{Zn(NO_3)_2} + \mathbf{2\mathrm{NO_2\uparrow}} + 2\mathrm{H_2O}
  \end{align}
  Passive metals ($\mathrm{Fe, Cr, Al}$) do not dissolve in concentrated $\mathrm{HNO_3}$ due to formation of an impervious surface oxide film.
  \item \textbf{Oxoacids of Phosphorus:}
  \begin{itemize}
    \item \textbf{Hypophosphorous acid ($\mathrm{H_3PO_2}$):} Monobasic ($1\ \mathrm{P-OH}$), contains $2\ \mathrm{P-H}$ bonds $\implies$ \textbf{powerful reducing agent} (reduces $\mathrm{AgNO_3 \to Ag\downarrow}$).
    \item \textbf{Phosphorous acid ($\mathrm{H_3PO_3}$):} Dibasic ($2\ \mathrm{P-OH}$), contains $1\ \mathrm{P-H}$ bond $\implies$ reducing agent. Disproportionates on heating: $4\mathrm{H_3PO_3 \xrightarrow{\Delta} 3H_3PO_4 + PH_3\uparrow}$.
    \item \textbf{Orthophosphoric acid ($\mathrm{H_3PO_4}$):} Tribasic ($3\ \mathrm{P-OH}$), zero $\mathrm{P-H}$ bonds $\implies$ non-reducing.
    \item \textbf{Pyrophosphoric acid ($\mathrm{H_4P_2O_7}$):} Tetrabasic ($4\ \mathrm{P-OH}$), contains $1\ \mathrm{P-O-P}$ bridge.
  \end{itemize}
\end{itemize}

\subsection{Group 16: Oxygen Family}
\begin{itemize}
  \item \textbf{Ozone ($\mathrm{O_3}$):} Prepared by passing silent electric discharge through pure dry dioxygen ($3\mathrm{O_2 \rightleftharpoons 2O_3}, \Delta H^\circ = +142\,\si{kJ.mol^{-1}}$). Quantitative estimation by adding neutral buffered $\mathrm{KI}$ ($\text{pH } 9.2$, borate buffer) and titrating liberated $\mathrm{I_2}$ against standard sodium thiosulfate:
  \begin{equation}
    2\mathrm{I^-} + \mathrm{H_2O} + \mathrm{O_3} \longrightarrow 2\mathrm{OH^-} + \mathrm{I_2} + \mathrm{O_2}
  \end{equation}
  Depletion by nitric oxide: $\mathrm{NO + O_3 \to NO_2 + O_2}$.
  \item \textbf{Allotropes of Sulfur:}
  \begin{itemize}
    \item \textbf{$\alpha$-Rhombic Sulfur:} Stable below $369\,\si{K}$ ($95.6^\circ\text{C}$), yellow, puckered crown $S_8$ ring.
    \item \textbf{$\beta$-Monoclinic Sulfur:} Stable above $369\,\si{K}$, needle-shaped, crown $S_8$ ring.
    \item \textbf{Transition Temperature ($369\,\si{K}$):} Both $\alpha$ and $\beta$ allotropes coexist in dynamic equilibrium.
    \item \textbf{Vapour Phase Paramagnetism:} Above $1000\,\si{K}$, sulfur exists predominantly as diatomic $\mathbf{\mathrm{S_2}}$ molecules which are \textbf{paramagnetic} due to two unpaired electrons in antibonding $\pi^*$ orbitals (identical to $\mathrm{O_2}$).
  \end{itemize}
  \item \textbf{Sulfuric Acid ($\mathrm{H_2SO_4}$ - Contact Process):}
  \begin{equation}
    2\mathrm{SO_2} + \mathrm{O_2} \xrightarrow[\mathrm{V_2O_5}, 720\,\si{K}]{2\,\si{bar}} 2\mathrm{SO_3} \quad (\Delta H^\circ = -196.6\,\si{kJ.mol^{-1}}), \quad \mathrm{SO_3} + \mathrm{H_2SO_4} \to \mathrm{H_2S_2O_7}\text{ (Oleum)}
  \end{equation}
  \item \textbf{Peroxo Oxoacids of Sulfur:}
  \begin{itemize}
    \item \textbf{Caro's Acid (Peroxomonosulfuric, $\mathrm{H_2SO_5}$):} Contains one $-O-O-$ peroxo bond. Oxidation state of $S$ is \textbf{$+6$} (NOT $+8$!).
    \item \textbf{Marshall's Acid (Peroxodisulfuric, $\mathrm{H_2S_2O_8}$):} Contains one $-O-O-$ peroxo bridge connecting two sulfur atoms. Oxidation state of $S$ is \textbf{$+6$}.
  \end{itemize}
\end{itemize}

\subsection{Group 17: Halogen Family}
\begin{itemize}
  \item \textbf{Bond Dissociation Enthalpy Anomaly:}
  \begin{equation}
    \mathbf{\mathrm{Cl_2 (242.6) > Br_2 (192.8) > F_2 (158.8) > I_2 (151.1)\,\si{kJ.mol^{-1}}}}
  \end{equation}
  Fluorine has an abnormally low bond dissociation energy due to intense electrostatic repulsion between non-bonding lone pairs on the tiny, closely spaced fluorine atoms.
  \item \textbf{Oxidizing Power Hierarchy:} $\mathbf{\mathrm{F_2 > Cl_2 > Br_2 > I_2}}$.
  Although $\Delta_{\text{eg}} H$ is more negative for Chlorine, Fluorine is the most powerful oxidizing agent because of:
  \begin{enumerate}
    \item Low bond dissociation enthalpy of $\mathrm{F_2}$.
    \item Extraordinarily high hydration enthalpy of the small $\mathrm{F^-}$ ion ($\Delta_{\text{hyd}} H = -515\,\si{kJ.mol^{-1}}$).
  \end{enumerate}
  \item \textbf{Hydrogen Halides ($\mathrm{HF, HCl, HBr, HI}$):}
  \begin{itemize}
    \item \textbf{Boiling Point:} $\mathbf{\mathrm{HF (293\,K) > HI (238\,K) > HBr (206\,K) > HCl (189\,K)}}$ ($\mathrm{HF}$ is liquid due to intermolecular H-bonding).
    \item \textbf{Acidic Strength:} $\mathrm{HF < HCl < HBr < HI}$ ($\mathrm{H-X}$ bond dissociation enthalpy decreases).
    \item \textbf{Reducing Power:} $\mathrm{HF < HCl < HBr < HI}$ ($\mathrm{HI}$ is the strongest reducing agent).
  \end{itemize}
  \item \textbf{Oxoacids of Halogens:}
  \begin{itemize}
    \item \textbf{Acidic Strength Order:} $\mathbf{\mathrm{HClO_4 > HClO_3 > HClO_2 > HClO}}$ (higher oxidation state of halogen increases conjugate base resonance stabilization via $d\pi-p\pi$ bonding).
    \item \textbf{Oxidizing Power Order:} $\mathbf{\mathrm{HClO > HClO_2 > HClO_3 > HClO_4}}$ (hypohalous acids are the most powerful oxidants due to ready oxygen atom liberation).
  \end{itemize}
\end{itemize}

\subsection{Group 18: Noble Gases}
Neil Bartlett (1962) observed that $\mathrm{PtF_6}$ oxidizes molecular dioxygen to $[\mathrm{O_2}]^+[\mathrm{PtF_6}]^-$. Because the first ionization enthalpy of Xenon ($1170\,\si{kJ.mol^{-1}}$) is nearly identical to that of $\mathrm{O_2}$ ($1175\,\si{kJ.mol^{-1}}$), he reacted Xenon with $\mathrm{PtF_6}$ to synthesize the first noble gas compound: $\mathrm{Xe^+[PtF_6]^-}$.

\begin{formulabox}[Xenon Fluorides: Preparation, Structures \& Hydrolysis Pathways]
\begin{align}
  \mathrm{Xe} + \mathrm{F_2} \xrightarrow[1\,\si{bar}]{673\,\si{K}} \mathrm{XeF_2} \quad (2:1\text{ ratio}) && \text{[Linear, } SN=5, 3\text{ lone pairs equatorial]} \\
  \mathrm{Xe} + 2\mathrm{F_2} \xrightarrow[7\,\si{bar}]{873\,\si{K}} \mathrm{XeF_4} \quad (1:5\text{ ratio}) && \text{[Square Planar, } SN=6, 2\text{ lone pairs axial]} \\
  \mathrm{Xe} + 3\mathrm{F_2} \xrightarrow[60-70\,\si{bar}]{573\,\si{K}} \mathrm{XeF_6} \quad (1:20\text{ ratio}) && \text{[Distorted Octahedral, } SN=7, 1\text{ lone pair]}
\end{align}
\textbf{Hydrolysis Pathways (Cardinal NEET Questions):}
\begin{align}
  2\mathrm{XeF_2} + 2\mathrm{H_2O} &\longrightarrow 2\mathrm{Xe} + 4\mathrm{HF} + \mathbf{\mathrm{O_2\uparrow}} \quad (\text{Redox hydrolysis}) \\
  6\mathrm{XeF_4} + 12\mathrm{H_2O} &\longrightarrow 4\mathrm{Xe} + 2\mathbf{\mathrm{XeO_3}} + 24\mathrm{HF} + 3\mathbf{\mathrm{O_2\uparrow}} \quad (\text{Disproportionation: } \mathrm{Xe^{IV}} \to \mathrm{Xe^0 + Xe^{VI}}) \\
  \mathrm{XeF_6} + 3\mathrm{H_2O} &\longrightarrow \mathbf{\mathrm{XeO_3}}\text{ (Explosive solid, pyramidal)} + 6\mathrm{HF} \quad (\text{Complete non-redox hydrolysis}) \\
  \mathrm{XeF_6} + \mathrm{H_2O} &\longrightarrow \mathbf{\mathrm{XeOF_4}}\text{ (Liquid, square pyramidal)} + 2\mathrm{HF} \quad (\text{Partial hydrolysis}) \\
  \mathrm{XeF_6} + 2\mathrm{H_2O} &\longrightarrow \mathbf{\mathrm{XeO_2F_2}}\text{ (Seesaw shape)} + 4\mathrm{HF} \quad (\text{Partial hydrolysis})
\end{align}
\textbf{Fluoride Ion Donor and Acceptor Reactions:}
\begin{equation}
  \mathrm{XeF_2} + \mathrm{PF_5} \longrightarrow [\mathrm{XeF}]^+[\mathrm{PF_6}]^-, \qquad \mathrm{XeF_6} + \mathrm{MF} \longrightarrow \mathrm{M}^+[\mathrm{XeF_7}]^- \quad (\mathrm{M} = \mathrm{Na, K, Rb, Cs})
\end{equation}
\end{formulabox}

\begin{ncertbox}[Direct NCERT Benchmark: $p$-Block Cardinal Diagnostic Extracts]
\begin{itemize}
  \item \textbf{Basicity vs Reducing Power in Group 15 Hydrides:} Basicity decreases ($\mathrm{NH_3 > PH_3 > AsH_3 > SbH_3 > BiH_3}$), whereas reducing power increases ($\mathrm{BiH_3 > SbH_3 > AsH_3 > PH_3 > NH_3}$). $\mathrm{BiH_3}$ is the strongest reducing hydride; $\mathrm{NH_3}$ is a mild reducing agent.
  \item \textbf{Phosphorus Oxoacid Reducing Power:} Arises exclusively from the presence of directly bonded $\mathrm{P-H}$ bonds:
  \begin{itemize}
    \item $\mathrm{H_3PO_2}$ (Hypophosphorous / Phosphinic): Has $2\ \mathrm{P-H}$ bonds $\implies$ most powerful reducing oxoacid.
    \item $\mathrm{H_3PO_3}$ (Phosphorous / Phosphonic): Has $1\ \mathrm{P-H}$ bond $\implies$ moderate reducing agent.
    \item $\mathrm{H_3PO_4}$ (Orthophosphoric): Has $0\ \mathrm{P-H}$ bonds $\implies$ strictly non-reducing.
  \end{itemize}
  \item \textbf{Paramagnetism of Diatomic Sulfur ($\mathrm{S_2}$):} In the vapour phase at elevated temperatures ($> 1000\,\si{K}$), sulfur exists as $\mathrm{S_2}$ molecules which possess two unpaired electrons in degenerate $\pi^*$ antibonding orbitals, exactly analogous to $\mathrm{O_2}$.
  \item \textbf{Commercial Preparation of Dinitrogen from Azides:} Ultra-pure nitrogen gas is produced by thermal decomposition of sodium or barium azide:
  \begin{equation*}
    \mathrm{2NaN_3 \xrightarrow{\Delta} 2Na + 3N_2\uparrow}, \qquad \mathrm{Ba(N_3)_2 \xrightarrow{\Delta} Ba + 3N_2\uparrow}
  \end{equation*}
\end{itemize}
\end{ncertbox}

% =========================================================================
% SECTION 6: D- AND F-BLOCK ELEMENTS
% =========================================================================
\section{Comparative Chemistry of $d$- and $f$-Block Elements (Unit 11)}

\subsection{General Trends of $3d$ Transition Series}
A transition element is defined as an element having an incompletely filled $d$-subshell in its ground state or in any of its common oxidation states. Zinc, Cadmium, and Mercury ($d^{10}$) are \textbf{non-transition elements} because they have completely filled $d$-orbitals in both ground and ionic states.
\begin{itemize}
  \item \textbf{Atomic Radii across $3d$ Series:} Decreases initially ($\mathrm{Sc}$ to $\mathrm{Cr}$), remains nearly constant ($\mathrm{Mn}$ to $\mathrm{Ni}$) because increasing nuclear charge is balanced by $d-d$ shielding, and slightly increases at the end ($\mathrm{Cu, Zn}$) due to interelectronic repulsion in filled $d^{10}$.
  \item \textbf{High Enthalpies of Atomization:} Reaches maximum at middle of series ($\mathrm{V, Cr, Mo, W}$) due to maximum unpaired electrons participating in strong interatomic metallic bonding.
  \item \textbf{Variable Oxidation States:} Arises because $(n-1)d$ and $ns$ electrons have comparable energies. $\mathrm{Sc}$ exhibits only $+3$; $\mathrm{Mn}$ exhibits the widest range from $+2$ to $+7$.
  \item \textbf{Standard Electrode Potentials ($E^\circ(M^{2+}/M)$):}
  \begin{equation}
    M(s) \xrightarrow{\Delta_{\text{sub}} H} M(g) \xrightarrow{\text{IE}_1 + \text{IE}_2} M^{2+}(g) \xrightarrow{\Delta_{\text{hyd}} H} M^{2+}(aq)
  \end{equation}
  $\mathbf{\mathrm{Cu}}$ has an anomalous \textbf{positive standard reduction potential ($+0.34\,\si{V}$)} because its high enthalpy of atomization and high ionization enthalpy are not compensated by the hydration enthalpy of $\mathrm{Cu^{2+}}$. Thus, Copper does not liberate $\mathrm{H_2}$ from non-oxidizing acids.
  \item \textbf{Disproportionation of Copper(I):}
  \begin{equation}
    2\mathrm{Cu^+(aq)} \longrightarrow \mathrm{Cu^{2+}(aq)} + \mathrm{Cu(s)} \quad (E^\circ_{\text{cell}} = +0.36\,\si{V})
  \end{equation}
  $\mathrm{Cu^{2+}}$ is far more stable in aqueous solution than $\mathrm{Cu^+}$ because the much higher hydration enthalpy of $\mathrm{Cu^{2+}}$ ($\Delta_{\text{hyd}} H \approx -2121\,\si{kJ.mol^{-1}}$) overcompensates for the second ionization enthalpy of Copper.
  \item \textbf{Spin-Only Magnetic Moment Formula:}
  \begin{equation}
    \mu_s = \sqrt{n(n + 2)}\,\si{BM}
  \end{equation}
  \begin{itemize}
    \item $n=1 \implies \mu = \sqrt{3} = 1.73\,\si{BM}$ ($\mathrm{Ti^{3+}, Cu^{2+}}$).
    \item $n=2 \implies \mu = \sqrt{8} = 2.83\,\si{BM}$ ($\mathrm{Ti^{2+}, Ni^{2+}}$).
    \item $n=3 \implies \mu = \sqrt{15} = 3.87\,\si{BM}$ ($\mathrm{V^{2+}, Cr^{3+}}$).
    \item $n=4 \implies \mu = \sqrt{24} = 4.90\,\si{BM}$ ($\mathrm{Cr^{2+}, Fe^{2+}}$).
    \item $n=5 \implies \mu = \sqrt{35} = 5.92\,\si{BM}$ ($\mathrm{Mn^{2+}, Fe^{3+}}$).
  \end{itemize}
\end{itemize}

\subsection{Master Transition Metal Compounds: $\mathrm{K_2Cr_2O_7}$ \& $\mathrm{KMnO_4}$}

\subsubsection{1. Potassium Dichromate ($\mathrm{K_2Cr_2O_7}$)}
\textbf{Industrial Preparation from Chromite Ore ($\mathrm{FeCr_2O_4}$):}
\begin{align}
  4\mathrm{FeCr_2O_4} + 8\mathrm{Na_2CO_3} + 7\mathrm{O_2} &\longrightarrow 8\mathrm{Na_2CrO_4\text{ (Yellow)}} + 2\mathrm{Fe_2O_3} + 8\mathrm{CO_2} \\
  2\mathrm{Na_2CrO_4} + 2\mathrm{H^+} &\longrightarrow \mathrm{Na_2Cr_2O_7\text{ (Orange)}} + 2\mathrm{Na^+} + \mathrm{H_2O} \\
  \mathrm{Na_2Cr_2O_7} + 2\mathrm{KCl} &\longrightarrow \mathbf{\mathrm{K_2Cr_2O_7\text{ (Orange crystals)}}} + 2\mathrm{NaCl}
\end{align}
$\mathrm{K_2Cr_2O_7}$ is preferred over $\mathrm{Na_2Cr_2O_7}$ as a primary volumetric standard because $\mathrm{Na_2Cr_2O_7}$ is deliquescent.
\textbf{pH-Dependent Chromate-Dichromate Interconversion:}
\begin{equation}
  \underset{\text{Yellow}}{\mathrm{2CrO_4^{2-}}} + 2\mathrm{H^+} \rightleftharpoons \underset{\text{Orange}}{\mathrm{Cr_2O_7^{2-}}} + \mathrm{H_2O} \quad (\text{Chromate in alkali; Dichromate in acid; Oxidation state stays }+6!)
\end{equation}
\textbf{Structures:}
\begin{itemize}
  \item Chromate ($\mathrm{CrO_4^{2-}}$): Tetrahedral ($T_d$), four equivalent $\mathrm{Cr-O}$ bonds ($166\,\si{pm}$) due to resonance.
  \item Dichromate ($\mathrm{Cr_2O_7^{2-}}$): Two tetrahedra sharing one oxygen corner; bridging $\mathrm{Cr-O-Cr}$ bond angle is \textbf{$126^\circ$}; bridging $\mathrm{Cr-O}$ bond length is $179\,\si{pm}$, while six terminal $\mathrm{Cr-O}$ bonds are shorter ($163\,\si{pm}$).
\end{itemize}
\textbf{Oxidizing Reactions in Acidic Medium ($n\text{-factor} = 6$):}
\begin{equation}
  \mathrm{Cr_2O_7^{2-}} + 14\mathrm{H^+} + 6e^- \longrightarrow 2\mathrm{Cr^{3+}} + 7\mathrm{H_2O} \quad (E^\circ = +1.33\,\si{V})
\end{equation}
\begin{align}
  6\mathrm{Fe^{2+}} &\longrightarrow 6\mathrm{Fe^{3+}} + 6e^- \\
  6\mathrm{I^-} &\longrightarrow 3\mathrm{I_2} + 6e^- \\
  3\mathrm{H_2S} &\longrightarrow 3\mathrm{S\downarrow} + 6\mathrm{H^+} + 6e^- \\
  3\mathrm{Sn^{2+}} &\longrightarrow 3\mathrm{Sn^{4+}} + 6e^-
\end{align}
\textbf{Chromyl Chloride Test:} Confirmation of chloride ions:
\begin{equation}
  \mathrm{K_2Cr_2O_7} + 4\mathrm{NaCl} + 6\mathrm{H_2SO_4(conc)} \xrightarrow{\Delta} 2\mathbf{\mathrm{CrO_2Cl_2\uparrow\text{ (Deep red vapors)}}} + 2\mathrm{KHSO_4} + 4\mathrm{NaHSO_4} + 3\mathrm{H_2O}
\end{equation}
Pass into $\mathrm{NaOH} \to$ yellow $\mathrm{Na_2CrO_4} \xrightarrow{\mathrm{Pb(CH_3COO)_2}}$ bright yellow precipitate of lead chromate ($\mathrm{PbCrO_4\downarrow}$).

\subsubsection{2. Potassium Permanganate ($\mathrm{KMnO_4}$)}
\textbf{Industrial Preparation from Pyrolusite Ore ($\mathrm{MnO_2}$):}
\begin{align}
  2\mathrm{MnO_2} + 4\mathrm{KOH} + \mathrm{O_2} &\xrightarrow{\text{fuse}} 2\mathbf{\mathrm{K_2MnO_4\text{ (Dark green)}}} + 2\mathrm{H_2O} \\
  3\mathrm{MnO_4^{2-}} + 4\mathrm{H^+} &\longrightarrow 2\mathbf{\mathrm{MnO_4^-\text{ (Purple)}}} + \mathrm{MnO_2} + 2\mathrm{H_2O} \quad (\text{Disproportionation})
\end{align}
\textbf{Manganate vs Permanganate Comparison:}
\begin{itemize}
  \item Manganate ($\mathrm{MnO_4^{2-}}$): Green, tetrahedral, $\mathrm{Mn(+6)}$, $3d^1 \implies$ \textbf{Paramagnetic} ($1$ unpaired electron).
  \item Permanganate ($\mathrm{MnO_4^-}$): Deep purple, tetrahedral, $\mathrm{Mn(+7)}$, $3d^0 \implies$ \textbf{Diamagnetic}.
  \textbf{Origin of Color:} Deep purple color is \textbf{NOT} due to $d-d$ transition (it is $d^0$!). Color arises from \textbf{Ligand-to-Metal Charge Transfer (LMCT)} where oxygen $2p$ electron density temporarily shifts into vacant manganese $3d$ orbitals.
\end{itemize}

\begin{formulabox}[Permanganate Redox Media \& Valence Factors ($n\text{-factor}$)]
\begin{enumerate}
  \item \textbf{In Acidic Medium ($n\text{-factor} = 5$):}
  \begin{equation}
    \mathrm{MnO_4^-} + 8\mathrm{H^+} + 5e^- \longrightarrow \mathrm{Mn^{2+}} + 4\mathrm{H_2O} \quad (E^\circ = +1.51\,\si{V})
  \end{equation}
  \begin{align}
    5\mathrm{Fe^{2+}} &\longrightarrow 5\mathrm{Fe^{3+}} + 5e^- \\
    5\mathrm{C_2O_4^{2-}} &\longrightarrow 10\mathrm{CO_2\uparrow} + 10e^- \quad (\text{titration heated to } 60-70^\circ\text{C}) \\
    10\mathrm{I^-} &\longrightarrow 5\mathrm{I_2} + 10e^- \\
    5\mathrm{NO_2^-} + 5\mathrm{H_2O} &\longrightarrow 5\mathrm{NO_3^-} + 10\mathrm{H^+} + 10e^- \\
    5\mathrm{SO_3^{2-}} + 5\mathrm{H_2O} &\longrightarrow 5\mathrm{SO_4^{2-}} + 10\mathrm{H^+} + 10e^-
  \end{align}
  \item \textbf{In Neutral / Faintly Alkaline Medium (Baeyer's Reagent, $n\text{-factor} = 3$):}
  \begin{equation}
    \mathrm{MnO_4^-} + 2\mathrm{H_2O} + 3e^- \longrightarrow \mathbf{\mathrm{MnO_2}} + 4\mathrm{OH^-} \quad (E^\circ = +1.23\,\si{V})
  \end{equation}
  \textbf{Cardinal NEET Reaction Differences:}
  \begin{itemize}
    \item Iodide is oxidized to \textbf{Iodate ($\mathrm{IO_3^-}$)}, NOT $\mathrm{I_2}$:
    \begin{equation}
      \mathrm{I^-} + 2\mathrm{MnO_4^-} + \mathrm{H_2O} \longrightarrow \mathbf{\mathrm{IO_3^-}} + 2\mathrm{MnO_2} + 2\mathrm{OH^-}
    \end{equation}
    \item Thiosulfate is oxidized to \textbf{Sulfate ($\mathrm{SO_4^{2-}}$)}:
    \begin{equation}
      8\mathrm{MnO_4^-} + 3\mathrm{S_2O_3^{2-}} + \mathrm{H_2O} \longrightarrow 8\mathrm{MnO_2} + 6\mathrm{SO_4^{2-}} + 2\mathrm{OH^-}
    \end{equation}
  \end{itemize}
  \item \textbf{In Strongly Alkaline Medium ($n\text{-factor} = 1$):}
  \begin{equation}
    \mathrm{MnO_4^-} + e^- \longrightarrow \mathrm{MnO_4^{2-}}
  \end{equation}
\end{enumerate}
\end{formulabox}

\begin{usecasebox}[High-Yield Problem Archetype: Permanganometric Volumetric Stoichiometry]
\textbf{Archetype Problem:} Determine the volume of $0.02\,\si{M}\ \mathrm{KMnO_4}$ required to completely oxidize:
(a) $25.0\,\si{mL}$ of $0.1\,\si{M}$ Mohr's salt ($\mathrm{FeSO_4\cdot(NH_4)_2SO_4\cdot 6H_2O}$), and
(b) $25.0\,\si{mL}$ of $0.1\,\si{M}$ Oxalic acid ($\mathrm{H_2C_2O_4\cdot 2H_2O}$) in acidic medium.
\begin{enumerate}
  \item \textbf{Equivalence Principle:} At the exact equivalence endpoint:
  \begin{equation*}
    N_1 V_1 = N_2 V_2 \iff (M_1 \times n_1) V_1 = (M_2 \times n_2) V_2
  \end{equation*}
  For $\mathrm{KMnO_4}$ in acidic medium, $n_1 = 5$ ($\mathrm{Mn^{VII} \to Mn^{II}}$).
  \item \textbf{Case (a): Titration with Mohr's Salt:}
  Mohr's salt acts as a single-electron reducing agent: $\mathrm{Fe^{2+} \to Fe^{3+} + e^-} \implies n_2 = 1$.
  \begin{equation*}
    (0.02 \times 5) \times V_1 = (0.1 \times 1) \times 25.0 \implies 0.10 \times V_1 = 2.50 \implies V_1 = \mathbf{25.0\,\si{mL}}
  \end{equation*}
  \item \textbf{Case (b): Titration with Oxalic Acid:}
  Oxalate ion undergoes a 2-electron oxidation: $\mathrm{C_2O_4^{2-} \to 2CO_2 + 2e^-} \implies n_2 = 2$.
  \begin{equation*}
    (0.02 \times 5) \times V_1 = (0.1 \times 2) \times 25.0 \implies 0.10 \times V_1 = 5.00 \implies V_1 = \mathbf{50.0\,\si{mL}}
  \end{equation*}
  \item \textbf{Practical Laboratory Condition:} The oxalic acid titration must be conducted at $60-70^\circ\text{C}$ because the initial reduction is slow. The reaction is \textbf{autocatalyzed} by the produced $\mathrm{Mn^{2+}}$ ions.
\end{enumerate}
\end{usecasebox}

\subsection{The $f$-Block Elements: Lanthanoids \& Actinoids}

\subsubsection{1. Lanthanoids ($4f$ Series: $\mathrm{Ce_{58}}$ to $\mathrm{Lu_{71}}$)}
\begin{itemize}
  \item \textbf{Electronic Configuration:} $[\mathrm{Xe}] 4f^{1-14} 5d^{0-1} 6s^2$.
  \item \textbf{Oxidation States:} Dominant and most stable state is \textbf{$+3$}.
  \begin{itemize}
    \item $+2$ exhibited by $\mathrm{Eu^{2+}}$ ($[\mathrm{Xe}] 4f^7$, half-filled) and $\mathrm{Yb^{2+}}$ ($[\mathrm{Xe}] 4f^{14}$, fully-filled) $\implies$ act as \textbf{strong reducing agents} reverting to $+3$.
    \item $+4$ exhibited by $\mathrm{Ce^{4+}}$ ($[\mathrm{Xe}] 4f^0$, empty shell) and $\mathrm{Tb^{4+}}$ ($[\mathrm{Xe}] 4f^7$) $\implies \mathbf{\mathrm{Ce^{4+}}}$ acts as a \textbf{powerful analytical oxidizing agent} ($E^\circ_{\mathrm{Ce^{4+}/Ce^{3+}}} = +1.74\,\si{V}$).
  \end{itemize}
  \item \textbf{Alloy: Mischmetal:} Consists of Lanthanoid metals ($\sim 95\%$), Iron ($\sim 5\%$), and traces of $\mathrm{S, C, Ca, Al}$. Used in cigarette lighter flints and tracer bullets.
\end{itemize}

\begin{ncertbox}[NCERT Hallmark: Lanthanoid Contraction \& Chemical Consequences]
The imperfect shielding of one $4f$ electron by another causes a progressive increase in effective nuclear charge ($Z_{\mathrm{eff}}$), resulting in a steady decrease in atomic and ionic radii from $\mathrm{Ce^{3+}}$ ($103\,\si{pm}$) to $\mathrm{Lu^{3+}}$ ($86\,\si{pm}$).
\textbf{Three Major Chemical Consequences for NEET:}
\begin{enumerate}
  \item \textbf{Covalent Character and Basicity of Hydroxides:} Size of $\mathrm{M^{3+}}$ decreases from $\mathrm{La^{3+}}$ to $\mathrm{Lu^{3+}}$. By Fajan's rules, covalent character of the $\mathrm{M-OH}$ bond increases; therefore, basic strength strictly decreases:
  \begin{equation*}
    \mathrm{La(OH)_3 > Ce(OH)_3 > \dots > Lu(OH)_3} \quad (\mathrm{La(OH)_3}\text{ is most basic, }\mathrm{Lu(OH)_3}\text{ is least basic})
  \end{equation*}
  \item \textbf{Pairing of 4d and 5d Transition Radii:} Due to the intervention of the 14 lanthanoid elements, elements of the 4d and 5d series show almost identical atomic and ionic radii:
  \begin{equation*}
    r(\mathrm{Zr}) \approx r(\mathrm{Hf}) \approx 160\,\si{pm}; \quad r(\mathrm{Nb}) \approx r(\mathrm{Ta}) \approx 146\,\si{pm}; \quad r(\mathrm{Mo}) \approx r(\mathrm{W})
  \end{equation*}
  Consequently, $\mathrm{Zr}$ and $\mathrm{Hf}$ exhibit identical chemical properties and occur together in nature, posing difficult separation challenges.
  \item \textbf{Chromatographic Separation:} Because radii vary by only $\sim 17\,\si{pm}$ across 14 elements, chemical separation is difficult, but their slight differences in complex formation constants and hydrated ionic radii permit successful separation via ion-exchange chromatography.
\end{enumerate}
\end{ncertbox}

\subsubsection{2. Actinoids ($5f$ Series: $\mathrm{Th_{90}}$ to $\mathrm{Lr_{103}}$)}
\begin{itemize}
  \item \textbf{Electronic Configuration:} $[\mathrm{Rn}] 5f^{1-14} 6d^{0-2} 7s^2$.
  \item \textbf{Oxidation States:} Much wider spectrum of oxidation states ($+3, +4, +5, +6, +7$) than Lanthanoids because the energy difference between $5f, 6d,$ and $7s$ subshells is exceedingly small.
  \item \textbf{Actinoid Contraction:} Steady decrease in ionic radii across the actinoid series. The contraction is \textbf{significantly greater} from element to element than the Lanthanoid contraction because $5f$ electrons exert even poorer shielding than $4f$ electrons.
  \item \textbf{Nuclear Instability:} All actinoids possess unstable nuclei; elements beyond Uranium ($Z > 92$) are synthetic \textbf{transuranic elements}.
\end{itemize}

\end{document}
